% !TEX root = ../print.tex
% Draft \S10 --- The corners: Matern, Student-t, symmetric stable
%
% Source: draft/spatial-hemigroup-scale-space.md, \S10 (lines 584-664).
%
% SPLIT (draft Proposition 10.2): the draft states the Matern family in six clauses, of
% which the exponent, the transforms and the moments are computations and the kernel density
% and its tail are classical Fourier pairs. The node is split so that the computations stay
% [T] and only the density and the tail are cited: prop:matern-exponent [T],
% prop:matern-density [A].
%
% CHANGED (prop:thorin-subclass clause 4): the draft states that the spatial Thorin measure
% is "the image of U_I under theta -> sqrt(2 theta)". It is TWICE that image. The factor is
% checked against the Matern member in the proof: with the causal Gamma data U_I = gamma
% delta_1 the image is gamma delta_{sqrt 2}, and (10.1) then returns half the right exponent.
% The draft should carry the factor 2.
%
% CHANGED (thm:matern clause 4): the draft's clause on rational stage transfer functions
% carries a check asking which rational functions are admissible; the general statement is
% prop:rational-increments. The node states the single-pole-pair case only, which is what the
% theorem uses, and refers the general case forward.

\chapter{The corners: \Matern{}, Student-t, symmetric stable}
\label{ch:corners}

\ref{thm:main-characterization} is a machine: one symmetric self-decomposable law in, one scale
space out. This chapter runs it on the three named families, says what each is, what cuts it
out of the cone and what it costs, and organises the three by the symmetric Thorin subclass, of
which the \Matern{} family is the single-atom member. Two general dials come first.

% draft: Proposition 10.1
% twin: hcs prop:moment-criterion (ledger A7 there) for clause (1); clause (2) has no causal
% node, the causal article stating the tail dichotomy in a remark.
% CHANGED (proof 2026-09-10, wave 4): clause (2)'s hypothesis is corrected. It read "if k
% vanishes beyond tau"; it now reads that tau is the LEAST such radius and that k is not
% identically zero. Both riders are needed and neither was decorative: at k == 0, the pure
% Gaussian, every tau satisfies the old hypothesis while the moment is finite at every alpha
% (Sato treats nu = 0 separately, Remark 26.3); and at a tau strictly larger than the radius the
% old statement asserted divergence on a range where its own first half asserts convergence.
% Ledger A14's entry already recorded the right reading -- "Sato's c is an infimum; the
% blueprint's tau is the same number" -- and the statement did not carry it. The Lean
% declaration Skeleton.moments_tails_bounded is repaired the same way, and the interface
% admitted on the trust boundary (SpatialLine.moments_tails_divergence) is the corrected
% divergence branch alone.
\begin{proposition}[moments and tails]\label{prop:moments-tails}
\uses{thm:main-characterization,lem:selfdecomposable-exponents}
\lean{SpatialLine.moments_tails_criterion, SpatialLine.moments_tails_variance,
Skeleton.moments_tails_bounded, SpatialLine.moments_tails_divergence,
Skeleton.moments_tails_heavy, SpatialLine.moments_tails_completely_monotone}\notready
Let $F$ be admissible with data $(a,k)$, and let $X_t$ be the displacement at canonical scale
$t$.
\begin{enumerate}
\item \emph{(Moments.)} For $n \ge 1$, $\mathbb{E}|X_t|^n < \infty$ if and only if
$\int_1^\infty x^{n-1}k(x)\,dx < \infty$. When finite, the variance is
$\mathbb{E}X_t^2 = t^2\bigl(2a + \int_0^\infty x\,k(x)\,dx\bigr)$, the two terms being the
jitter and the second moment of the catalogue.
\item \emph{(Tails.)} If $k \not\equiv 0$ and $\tau$ is the least radius beyond which $k$
vanishes, then $\mathbb{E}\exp(\alpha|X_t|\log|X_t|) < \infty$ for every $\alpha < 1/(\tau t)$
and for no $\alpha > 1/(\tau t)$: the tail is of order $e^{-|x|\log|x|/(\tau t)}$. If $k$
vanishes beyond no $\tau$, the tail is heavier than every such rate. In particular a completely monotone $k$,
being positive on all of $(0,\infty)$, puts the whole Thorin subclass on the heavy side, and no
kernel other than the Gaussians themselves has a Gaussian or lighter tail.
\end{enumerate}
\statusA\quad\ledger{A11}\quad\ledger{A13}\quad\ledger{A14}\quad\emph{(\textbf{Assignment.} \ledger{A13}
carries the moment criterion for infinitely divisible laws, in the submultiplicative-weight
form, and \ledger{A14} the two-sided tail statement for a \Levy{} law whose jump measure has
bounded support. What they do not carry: the translation into the folded profile --- the \Levy{}
measure of $X_t$ is the $t$-dilate of $k(x)x^{-1}dx$ folded, so
$\int_{|x|>1}|x|^n\,\nu_2(dx) = \int_1^\infty x^{n-1}k(x)\,dx$ up to the dilation --- and the
variance formula, which is the second difference $2\omega^{-2}(1 - \hat\mu_t(\omega))$ at the
origin evaluated on both sides; both are \textbf{[T]}, elementary, and written out below. The
variance clause in particular spends no ledger entry at all: the finiteness of
$\mathbb{E}X_t^2$ is a conclusion of its own argument and not an appeal to \ledger{A13}. The last
sentence of clause (2) spends one entry more, and \ledger{A11} is named above for it: that a
completely monotone profile is a mixture of exponentials is Bernstein's theorem, and it is how the
profile is seen to be positive at every $x > 0$; the comparison of the two regimes that follows is
\textbf{[T]}. (\texttt{SpatialLine.moments\_tails\_completely\_monotone} reaches
\texttt{bernstein\_completely\_monotone} through
\texttt{completely\_monotone\_support\_unbounded}. The status line did not name the entry, and the
annotation called the whole sentence \textbf{[T]}, until 2026-09-19; see the marker below.)
Of \ledger{A14} only the divergence branch is admitted in Lean,
and only at a lower bound for the radius, which is what the last sentence consumes; the
finiteness branch of clause (2) is cited and not admitted, no proof spending it.)}
\end{proposition}

% CHANGED (2026-09-19, article mirror of module B, R174): TWO CHANGES, neither to a clause.
% (a) THE STATUS LINE NAMES \ledger{A11}, and the annotation says where it is spent. The last
% sentence of clause (2) -- that a completely monotone profile is positive on all of (0,infinity),
% so that the whole Thorin subclass sits on the heavy side -- passes through Bernstein's theorem,
% in the prose proof ("a completely monotone k is a mixture of exponentials") and in the Lean
% (SpatialLine.moments_tails_completely_monotone reaches bernstein_completely_monotone through
% completely_monotone_support_unbounded, IsCompletelyMonotone being the derivative condition). The
% node named A13 and A14 only, and the annotation called that sentence [T]. This is a DISCLOSURE,
% in the safe direction: the article names one more cited fact and claims no more. No statement,
% tag, status or declaration changes, and A11 is admitted already (as the existence equivalence,
% at prop:thorin-subclass), so no trust-boundary name is added.
% (b) The proof's use of the THORIN MEASURE is replaced by an anonymous representing measure (the
% register pass's item 22: "U" appeared here, in the chapter's first proof, long before
% prop:thorin-subclass introduces it). Bernstein's theorem is named where it is used.
% CHANGED (2026-09-11, draft sync, R103): a "proof of record 2026-09-09, wave 3" block stood
% here describing the equivalence of clauses (1) and (2), sigma-finiteness of U, the averaging
% step and uniqueness of U -- which is prop:thorin-subclass's proof, not this node's (this node
% has no such clauses and no U). It is moved, verbatim, to above prop:thorin-subclass's proof.
% CHANGED (proof of record 2026-09-10, wave 5): the VARIANCE paragraph is rewritten to the
% machine-checked route. The printed route differentiates eq:levy-khintchine twice under the
% integral sign; the obligation is a second difference of the transform at the origin, which
% needs one dominated convergence rather than two differentiations and which delivers the
% finiteness of E X_t^2 rather than assuming it (Fatou at the origin). The printed route is
% kept as a marked remark. The consequence for the axiom line is recorded in the annotation:
% SpatialLine.moments_tails_variance prints Lean core, so ledger A13 is not on the variance
% clause's path.
\begin{proof}[Proof of the parts held \textbf{[T]}]
The \Levy{} measure of $X_t$ is the $t$-dilate of $\nu = k(x)x^{-1}dx$, folded to the two-sided
measure $\nu_2$; so
$\int_{|x|>1}|x|^n\,\nu_2(dx) = \int_{1/t}^\infty (tx)^n\,k(x)\,x^{-1}dx$, which is finite
exactly when $\int_1^\infty x^{n-1}k(x)\,dx$ is, the integrand being locally bounded away from
the origin. Applying A13 with the submultiplicative weight $\max(|x|,1)^n$ gives clause (1)'s
criterion: that weight agrees with $|x|^n$ on $\{|x| > 1\}$, the only region the criterion
integrates over, and has a finite integral against the law exactly when $|x|^n$ does, since
$|x|^n \le \max(|x|,1)^n \le |x|^n + 1$; the weight $|x|^n$ itself is not submultiplicative.
% CHANGED (2026-09-11, draft sync, R104): read "Applying A13 with the weight |x|^n". That weight
% is not submultiplicative in Sato's Def. 25.2, so Thm. 25.3 does not apply to it; the entry's
% weight is (|x| v 1)^n since R49 corrected blueprint/AXIOMS.md. Textual: the two moments are
% finite together and the weights agree where the criterion integrates. Mirrored in the draft.

For the variance, what has to be evaluated is a second difference of the transform at the
origin and not a second derivative. Write $C = t^2(2a + \int_0^\infty x\,k(x)\,dx)$, finite by
hypothesis. The half-angle formula $1 - \cos u = 2\sin^2(u/2)$ gives, for every real $c$ and
every $\omega \ne 0$,
\[
  \bigl(1 - \cos c\omega\bigr)\frac{2}{\omega^2}
   = c^2\Bigl(\frac{\sin(c\omega/2)}{c\omega/2}\Bigr)^2 \le c^2 ,
\]
an identity rather than an estimate, so the same line supplies the pointwise limit $c^2$ as
$\omega \to 0$, a dominating function, and the inequality $1 - \cos u \le u^2/2$. On the
exponent side, reading it at $c = tx$ inside \ref{eq:sd-profile} and dominating by $t^2xk(x)$
gives $2\omega^{-2}F(t\omega) \to C$; hence $F(t\omega) \to 0$, and since
$ye^{-y} \le 1 - e^{-y} \le y$ for $y \ge 0$, the same limit holds for
$2\omega^{-2}(1 - e^{-F(t\omega)})$. On the measure side, $\mu_t$ being a probability measure,
$1 - \hat\mu_t(\omega) = \int(1 - \cos\omega x)\,\mu_t(dx)$; Fatou's lemma along a sequence
$\omega_m \to 0$ turns the exponent-side limit into $\int x^2\,\mu_t(dx) \le C$, so the second
moment is finite --- a conclusion of the argument, not a hypothesis of it --- and dominated
convergence with the dominating function $x^2$ then gives
$2\omega^{-2}(1 - \hat\mu_t(\omega)) \to \int x^2\,\mu_t(dx)$. The two limits are limits of the
same function, so $\mathbb{E}X_t^2 = C$. The mean vanishes because $|x| \le 1 + x^2$ makes $x$
integrable against $\mu_t$ and $\mu_t$ is reflection-symmetric.

\emph{Remark, not machine-checked.} A second route to the variance evaluates
$\mathbb{E}X_t^2 = -\partial_\omega^2 e^{-F(t\omega)}|_{\omega = 0} = t^2F''(0)$ and
differentiates \ref{eq:levy-khintchine} twice under the integral sign to obtain
$F''(0) = 2a + \int_{(0,\infty)}x^2\,\nu(dx)$. It reaches the same constant and costs a
dominating function for each of the two differentiations, where the second difference costs one
and delivers the finiteness of $\mathbb{E}X_t^2$ along the way.

For the last
sentence of (2), a completely monotone $k$ is a mixture of exponentials by Bernstein's theorem
(A11): $k(x) = \int e^{-\theta x}\,V(d\theta)$ for a positive measure $V$, which is not the zero
measure because $k \not\equiv 0$. Some window $[0,\theta_0]$ then carries positive mass, and on it
$e^{-\theta x} \ge e^{-\theta_0x}$, so $k(x) \ge e^{-\theta_0x}\,V([0,\theta_0]) > 0$ at every
$x > 0$; hence $k$ vanishes beyond no $\tau$, and the heavy branch of A14 applies. Finally, if $k \equiv 0$
then $F = a\omega^2$ and the kernel is Gaussian, while if $k \not\equiv 0$ then $F$ has a
nonzero jump catalogue and A14's tail floor forbids a Gaussian tail: fixing $x_0 > 0$ with
$k(x_0) \ne 0$, the catalogue reaches the radius $x_0$, so the floor holds at every
$\alpha' > 1/(x_0t)$, and one such $\alpha'$ suffices, since
$\alpha'|x|\log|x| \le \alpha x^2$ once $|x| \ge (2\alpha'/\alpha)^2$ --- write
$\log|x| = 2\log\sqrt{|x|} \le 2\sqrt{|x|}$ --- while the law is a probability measure, so the
divergence cannot come from the bounded part. No case distinction between a bounded and an
unbounded catalogue is needed at this step, one radius the catalogue reaches being enough.
\end{proof}

The two dials are turned by one object, the tail of $k$: polynomial moment weight by its decay,
super-exponential tail order by its support. The $\Cin$ rays and their superpositions, with
bounded support, are the best-localized members of the cone and have every moment; the price is
the echo structure of \ref{lem:cin-delay-equation}. Completely monotone profiles buy echo-free
kernels realizable by first-order sections (Chapter~\ref{ch:implementation}) at the price of
tails no lighter than exponential.

% draft: Proposition 10.2, clauses (1), (2), (4) --- SPLIT, see the header note.
% twin: hcs prop:gamma-family + prop:gamma-kernels + prop:gamma-moments
% (Hemigroup.SelfDecomposableExponent.gammaExponent_toRealExponent,
%  ...gammaExponent_laplace_kernel, ...gammaExponent_moments)
% CHANGED (proof of record 2026-09-10, wave 5): clause (3)'s proof is rewritten to the
% machine-checked route, and the \uses edge to prop:bridge-families is REMOVED with it. The
% printed proof imported the Gamma mixture from that node; the mixing law is a Gamma law with a
% known Laplace transform and is CONSTRUCTED here in three lines, after which
% prop:fourier-uniqueness identifies the two measures. So clause (3) does not depend on
% chapter 9 at all, and the edge said it did. The identification of the Matern family as the
% image of the causal Gamma family is still prop:bridge-families(2)'s content and that node
% still cites this one; the direction of the dependency is what changed. The variance route
% through prop:moments-tails(1) survives as a marked remark: it is machine-checked, it is
% independent, and it does not reach the general even moment.
% CHANGED (author's decision 2026-09-10): SpatialLine.matern_moments_integrable is DROPPED from
% the \lean tag. It is the A13 route of the remark at the end of the proof -- a second and
% independent part-proof of the finiteness conjunct, machine-checked but spending a
% trust-boundary name -- and naming it beside matern_moments made a [T] node's tag list an
% axiom-spending declaration, which is exactly the shape a blind fidelity reviewer should query.
% The declaration stays in SpatialLine/Moments.lean as a corollary, with a docstring line
% saying it is no longer a node declaration; A13's ledger entry still lists it as a consumer,
% which remains true, and the entry's own [A] node is prop:moments-tails, not this one. The
% node's four remaining declarations print Lean core, so the tag and the \statusT now agree.
\begin{proposition}[the \Matern{} family: exponent, transforms, moments]\label{prop:matern-exponent}
\uses{lem:selfdecomposable-exponents,lem:profile-integrability,prop:moments-tails,prop:fourier-uniqueness,thm:main-characterization}
\lean{SpatialLine.matern_exponent, SpatialLine.matern_transforms,
SpatialLine.matern_moments, SpatialLine.matern_gamma_mixture}\leanok
\notes{gamma-matern-corner}
For $\gamma > 0$:
\begin{enumerate}
\item \emph{(Exponent.)} $F(\omega) = \gamma\log(1 + \omega^2)$ is admissible, with Gaussian
coefficient $a = 0$ and folded profile $k(x) = 2\gamma e^{-x}$; the profile at range $\theta$ is
$k(x) = 2\gamma e^{-x/\theta}$, with exponent $\gamma\log(1 + \theta^2\omega^2)$.
\item \emph{(Transforms.)} $e^{-F(t\omega)} = (1 + t^2\omega^2)^{-\gamma}$, and the increments
have transfer functions $\bigl((1 + s^2\omega^2)/(1 + t^2\omega^2)\bigr)^{\gamma}$, rational in
$\omega$ for integer $\gamma$.
\item \emph{(Moments.)} Every moment is finite: $X_t \seq \BM_T$ with $T$ Gamma-distributed of
shape $\gamma$ and scale $2t^2$, so
$\mathbb{E}X_t^{2n} = (2n-1)!!\,(2t^2)^n\,\Gamma(\gamma+n)/\Gamma(\gamma)$, and in particular
$\mathbb{E}X_t^2 = 2\gamma t^2$.
\end{enumerate}
\statusT\quad\emph{(Every clause is machine-checked, and the node's footprint is Lean core:
finiteness of every moment follows from the even moments by $|x|^n \le 1 + x^{2n}$, and the
mixture is constructed rather than cited. \textbf{Changed 2026-09-10 (author's decision).} The
Lean tag used to name a fifth declaration, \texttt{SpatialLine.matern\_moments\_integrable},
which proves the finiteness conjunct a second time through \ref{prop:moments-tails}(1) and
spends ledger A13 in doing so. That is the remark below: a genuine and independent part-proof,
but not the node's grounding, and a $[T]$ node whose tag names an axiom-spending declaration
reads as though it were. The declaration stays in the library as a corollary; the tag now names
only what the node rests on, so what the tag lists and what the node's status claims agree.)}
\end{proposition}

% CHANGED (proof of record 2026-09-09): clause (1)'s proof is rewritten to the machine-checked
% route, which is shorter than the printed one in two places and longer in one. Longer: the
% domination that makes the differentiation legitimate is now written out, because it is the
% whole content of the step and because it is uniform in omega, which is what lets the identity
% be proved on all of R at once rather than on a neighbourhood. Shorter: (i) the two
% integrability conditions are the SAME two dominating functions read on (0,1) and (1,infinity),
% so lem:profile-integrability is not spent -- that lemma converts the \Levy{} condition into these
% two, and here the two are what is checked directly; (ii) lem:selfdecomposable-exponents is not
% spent either. Exhibiting the pair (a,k) of eq:sd-profile IS admissibility in the sense of
% thm:main-characterization, so the appeal was a step the statement does not need; and the range
% theta comes out of the computation itself, so it is not obtained from that lemma's dilation
% identity. The uses edges are kept: clauses (2) and (3) still consume those nodes.
\begin{proof}
\leanok
(1) Fix $\theta > 0$ and write $p = 1/\theta$. The integrand
$(1 - \cos\omega x)\,e^{-px}x^{-1}$ is nonnegative on $(0,\infty)$ and, since
$1 - \cos u \le u^2/2$, is bounded there by $\tfrac12\omega^2\,x\,e^{-px}$, so it is integrable
for every $\omega$; its derivative in $\omega$ is $\sin(\omega x)\,e^{-px}$, whose modulus is at
most $e^{-px}$ \emph{uniformly in} $\omega$. Differentiation under the integral sign is
therefore legitimate at every $\omega$ with one and the same dominating function, and gives
\[
  \frac{d}{d\omega}\int_0^\infty (1 - \cos\omega x)\,e^{-px}\,\frac{dx}{x}
   = \int_0^\infty \sin(\omega x)\,e^{-px}\,dx = \frac{\omega}{p^2+\omega^2},
\]
which is also the derivative of $\tfrac12\log(1 + \omega^2/p^2)$. The two functions of $\omega$
therefore differ by a constant, and both vanish at $\omega = 0$, so
\[
  \int_0^\infty (1 - \cos\omega x)\,e^{-x/\theta}\,\frac{dx}{x}
   = \tfrac12\log\bigl(1 + \theta^2\omega^2\bigr).
\]
Hence $F(\omega) = \gamma\log(1+\theta^2\omega^2)$ is of the form \ref{eq:sd-profile} with
$a = 0$ and $k(x) = 2\gamma e^{-x/\theta}$, which is nonnegative and nonincreasing on
$(0,\infty)$. Its two integrability conditions are the two dominating functions above read on
the two halves of the range: $\int_0^1 x\,k(x)\,dx \le 2\gamma$ and
$\int_1^\infty k(x)\,x^{-1}dx \le 2\gamma\int_1^\infty e^{-x/\theta}dx < \infty$. Exhibiting the
pair $(0, 2\gamma e^{-\cdot/\theta})$ is what admissibility asks for, so the argument stops
there; the range $\theta$ is carried by the computation and needs no separate dilation step.

(2) is (1) read through the similarity form of \ref{thm:main-characterization}: the increment
transfer function is the ratio of the endpoint transforms, the denominators being nonvanishing.

(3) Let $\Gamma_{\gamma,2t^2}$ be the Gamma law of shape $\gamma$ and scale $2t^2$, that is the
law on $(0,\infty)$ with density $u^{\gamma-1}e^{-u/2t^2}/(\Gamma(\gamma)(2t^2)^{\gamma})$. Its
Laplace transform is $\int_0^\infty e^{-\sigma u}\,\Gamma_{\gamma,2t^2}(du) =
(1 + 2t^2\sigma)^{-\gamma}$ for $\sigma \ge 0$, one Gamma integral; so the Gaussian mixture
$\int_0^\infty g_u\,\Gamma_{\gamma,2t^2}(du)$ --- the law of $\BM_T$ with $T$ so distributed ---
has cosine transform $(1 + 2t^2\cdot\omega^2/2)^{-\gamma} = (1 + t^2\omega^2)^{-\gamma}$, which
by (2) is the transform of the kernel at scale $t$. Both are symmetric probability measures, so
\ref{prop:fourier-uniqueness} identifies them and $X_t \seq \BM_T$. Conditioning on $T$,
$\mathbb{E}X_t^{2n} = \mathbb{E}[T^n]\,\mathbb{E}[\BM_1^{2n}] = (2n-1)!!\;\mathbb{E}T^n$ with
$\mathbb{E}[\BM_v^{2n}] = (2n-1)!!\,v^n$ and
$\mathbb{E}T^n = (2t^2)^n\Gamma(\gamma+n)/\Gamma(\gamma)$; the two are the same Gamma integral
read at half-integer and at integer argument. Every moment is then finite, since
$|x|^n \le 1 + x^{2n}$, and $n = 1$ with $\Gamma(\gamma+1) = \gamma\Gamma(\gamma)$ gives
$\mathbb{E}X_t^2 = 2\gamma t^2$.

\emph{Remark.} Finiteness of every moment also follows from \ref{prop:moments-tails}(1), the
profile decaying exponentially, and the variance from the variance formula there:
$\mathbb{E}X_t^2 = t^2\int_0^\infty x\,2\gamma e^{-x}dx = 2\gamma t^2$. That route is machine-
checked too and is independent of this one; it does not reach the general even moment.
\end{proof}

% draft: Proposition 10.2, clauses (3), (5), (6) --- SPLIT, see the header note.
% SKELETON NOTE (phase B, 2026-09-09): Mathlib (v4.31.0) carries no modified Bessel function
% and no error function, so SpatialLine/Corners.lean DEFINES besselK by the integral
% representation DLMF (10.32.9) -- which is what this entry's own sources define it by -- and
% erfc by its defining integral. The Lean statements of this node, of prop:student-t(2), of
% thm:joint-locality(2) and of prop:thorin-subclass(5) are against those definitions. The
% consequence to keep in view is that A16 and A19 become statements about that integral, with
% no Mathlib API behind the name; SKELETON.md question Q10 asks the review to confirm the
% primitive. The gamma = 1 special value needs K_{1/2}(z) = sqrt(pi/2z) e^{-z} and is priced L
% for that reason, while the gamma = 1/2 value is an S computation from the definitions.
% CHANGED (review 2026-09-09, R20): the power of t in the tail was WRONG and is corrected from
% t^{-gamma-1/2} to t^{-gamma} (the draft, \S10, carries the same error). From the scaling
% phi_t = t^{-1}phi_1(./t) and phi_1(y) ~ c|y|^{gamma-1}e^{-|y|} one gets t^{-1}(|x|/t)^{gamma-1}
% = t^{-gamma}|x|^{gamma-1}; the same follows from K_nu(z) ~ sqrt(pi/2z)e^{-z} in the displayed
% density, and the check at gamma = 1 settles it, the Laplace kernel being e^{-|x|/t}/(2t) with
% t-power -1 and not -3/2. The Lean statement (Skeleton.matern_density_tail) had the constant
% quantified INSIDE the scope of t, so it absorbed the error and asserted nothing about t; it
% now quantifies the constant first, which is what makes the exponent a claim.
% CHANGED (proof 2026-09-09, wave 2): Skeleton.matern_density_gaussian_limit was FALSE as typed
% and gained the hypothesis that the kernels are SYMMETRIC. It is the review's R7 defect at a
% third site. The declaration quantified over probability measures constrained only by their
% COSINE transform -- which determines a measure's symmetrisation and nothing more -- and
% concluded weak convergence, a statement about the whole measure. Counterexample: add to the
% Matern law nu_n the odd signed measure of density (1/2) nu_n(x) sgn(x); the sum is a
% probability measure with the same cosine transform, and for the bounded continuous odd
% g = arctan its integral tends to (1/2)E|arctan N| > 0 while the Gaussian limit gives 0. The
% repair is R7's own -- symmetry of the kernels, which (A3) gives for every family here and
% which prop:two-members already assumes -- and it narrows nothing. The statement of record
% below is unaffected: the blueprint says "the kernel of the Matern family", and that family's
% kernels are symmetric; what was missing was in the LEAN quantification, not in the prose.
% CHANGED (2026-09-14, R156; a NARROWING of the naming sentence): the density is called the
% Matern kernel of smoothness gamma - 1/2 only where it is one, gamma > 1/2. The displayed
% function is a probability density for every gamma > 0, but a Matern COVARIANCE function is
% bounded at the origin, which needs smoothness gamma - 1/2 > 0; at gamma = 1/2 it is K_0, which
% diverges logarithmically. The variance-gamma identification, which is what the ledger entry
% A15 carries and what the later nodes read, holds on the whole range and is stated so. Same
% finding as R156 at prop:two-members (the external review's finding C); mirrored in the draft.
\begin{proposition}[the \Matern{} family: the kernel and its tail]\label{prop:matern-density}
\uses{prop:matern-exponent,prop:kernel-regularity,prop:levy-continuity,prop:moments-tails}
\lean{Skeleton.matern_density, SpatialLine.matern_density_special, Skeleton.matern_density_tail,
SpatialLine.matern_density_gaussian_limit}\notready
\notes{gamma-matern-corner}
For $\gamma > 0$ and $t > 0$ the kernel of the \Matern{} family at canonical scale $t$ has the
density
\[
  \phi_t(x) = \frac{1}{\sqrt\pi\,\Gamma(\gamma)\,t}
   \Bigl(\frac{|x|}{2t}\Bigr)^{\gamma - \frac12}
   K_{\gamma-\frac12}\Bigl(\frac{|x|}{t}\Bigr),
\]
for $\gamma > \tfrac12$ the \Matern{} kernel of smoothness $\gamma - \tfrac12$ and range $t$
normalised to unit mass; for every $\gamma > 0$ these are the symmetric variance-gamma laws. Here
$\gamma = 1$ gives the Laplace kernel $e^{-|x|/t}/(2t)$ at every $x \ne 0$, and
$\gamma = \tfrac12$ gives $K_0(|x|/t)/(\pi t)$ at every $x$. The origin is excluded from the first
of those two values because the display reads there as the product of $(|x|/2t)^{1/2} = 0$ with
$K_{1/2}(0)$ and does not return $(2t)^{-1}$. The difference is on a null set and leaves the
kernel, a density, unchanged.
% CHANGED (2026-09-19, article mirror of module B, R174): the two special values and the null-set
% exclusion are two sentences where they were one with an em-dash aside (the register pass's
% item 13). The gamma = 1/2 value gains the words "at every x", which is what the R57 comment below
% records of it and which the contrast with the gamma = 1 value now makes explicit. No claim
% changes: the same two values, the same exclusion, the same reason.
% CHANGED (2026-09-11, draft sync, R57): the gamma = 1 special value was stated without a domain.
% SpatialLine.matern_density_special proves it for x /= 0 only: at x = 0 the closed form carries
% the factor (|0|/2t)^{1/2} = 0 and the Laplace density is (2t)^{-1}. Narrowed to x /= 0, a
% Lebesgue-null exception that leaves the kernel, a density, unchanged; the gamma = 1/2 value has
% no exclusion (both sides read K_0(0)/(pi t) at the origin). Mirrored at draft Proposition 10.2(3).
% CHANGED (2026-09-15, module B drafting, R164): the words "As |x| -> infinity," opened the tail
% sentence and had been swept into the R57 comment line above by an editing slip, so the
% statement of record read "phi_t(x) ~ ..." with no limit named; the checker strips comments, so
% the printed statement lacked the words while the draft and the Lean (Skeleton.matern_density_tail,
% a statement at infinity) had them. Restored on the statement's own line. Found when the paper
% transcribed the node verbatim.
As $|x| \to \infty$,
$\phi_t(x) \sim c_\gamma\,t^{-\gamma}|x|^{\gamma-1}e^{-|x|/t}$: exponential tails, the
heavy side of \ref{prop:moments-tails}(2), as for every completely monotone profile. As
$\gamma \to \infty$ with $2\gamma t^2 = v$ held fixed, $\phi_t$ converges weakly to the
Gaussian of variance $v$.
\statusA\quad\ledger{A15}\quad\ledger{A16}\quad\emph{(\textbf{Assignment.} \ledger{A15} carries
the Fourier pair, that the displayed density has transform $(1+t^2\omega^2)^{-\gamma}$, together
with the identification of the family as the symmetric variance-gamma laws; \ledger{A16} carries
the asymptotics of the modified Bessel function that give the tail. What they do not carry: the
two special values, which are the displayed formula at $\gamma = 1$ and $\gamma = \tfrac12$;
that the kernel is a density at all, which is \ref{prop:kernel-regularity}; and the Gaussian
limit, which is \textbf{[T]} from $(1 + v\omega^2/2\gamma)^{-\gamma} \to e^{-v\omega^2/2}$ and
\ref{prop:levy-continuity}.)}
\end{proposition}

% draft: Theorem 10.3
% twin: hcs has no single node for this; the causal article states the Gamma corner across
% prop:gamma-family, prop:gamma-kernels, prop:gamma-density and prop:thorin-subclass.
% CHANGED (review 2026-09-09, Q11): the closing CONSEQUENCE -- that the Matern family is the
% unique admissible family with every moment finite, exponential tails and an exact first-order
% realization -- has MOVED OUT of this theorem into cor:matern-realization below, together with
% the \uses edge to prop:matern-cascade and the closing paragraph of the proof. It rested on a
% Chapter 14 node with no Lean scheduled, so it could not be typed, and a \lean tag that covers
% only some of a node's clauses is the exception to the fidelity rule this split removes: the
% tag below now names declarations for the equivalence, which is all the theorem asserts.
\begin{theorem}[the \Matern{} theorem]\label{thm:matern}
\uses{prop:matern-exponent,prop:matern-density,prop:thorin-subclass,prop:fourier-toolbox}
\lean{SpatialLine.matern_theorem, SpatialLine.matern_thorin_atom, SpatialLine.matern_kernels,
SpatialLine.matern_rational, SpatialLine.matern_rational_polynomial,
SpatialLine.matern_gauge_profile, SpatialLine.matern_gauge}\leanok
\notes{gamma-matern-corner}
% CHANGED (2026-09-19, external review of module B, R170): two faults of the printed statement
% against its declaration. (a) Clause (2) lacked the hypothesis of no Gaussian part:
% a omega^2 + gamma log(1 + theta^2 omega^2) has a one-point Thorin measure and satisfies neither
% (1) nor (3). SpatialLine.matern_theorem has `P.a = 0` and `thorinExponent 0` in that clause; the
% print had dropped it (the fault of R167 again). (b) "The following are equivalent" covered
% clause (4), which is an equivalence for a positive integer gamma only and is a separate
% conjunct over m : N in the declaration; the opening sentence now says which clauses are
% equivalent for which gamma, and "a single pole pair" is said of s < t, the stage at s = t being
% the identity. Found by the referee. Safe direction; no Lean change.
% CHANGED 2026-09-19 (R176, second external review round, referee B finding 2): THE FAMILY IS IN
% ITS CANONICAL GAUGE, a hypothesis added. Clauses (3) and (4) are about the family's own kernels
% and stages at the parameter t, so they are gauge-dependent: the referee's Phi'_{s,t} :=
% Phi_{s^2,t^2} of a Matern family satisfies the axioms and (ND), has clause (1)'s profile, and
% has stage transfer ((1 + theta^2 chi(s)^2 omega^2)/(1 + theta^2 chi(t)^2 omega^2))^gamma with
% chi =/= id, so clause (4) as printed failed for it. The Lean had the hypothesis all along:
% SpatialLine.matern_theorem takes hcos : fourierCos (mu 0 t) omega = exp(-P.exponent (t * omega)),
% which IS the canonical gauge. So the print is brought to the declaration; the same hypothesis is
% what prop:bridge-families already states where it needs it. Safe direction (a hypothesis added);
% no Lean change. cor:matern-realization gets it too, for the same reason.
For an admissible family \emph{in its canonical gauge} and $\gamma, \theta > 0$, clauses (1), (2)
and (3) are equivalent, and
when $\gamma$ is a positive integer they are equivalent to clause (4).
\begin{enumerate}
\item The profile is a single exponential, $k(x) = 2\gamma e^{-x/\theta}$, and $a = 0$.
\item The family lies in the symmetric Thorin subclass of \ref{prop:thorin-subclass} with
Gaussian coefficient $0$ and the one-point Thorin measure $2\gamma\,\delta_{1/\theta}$.
\item The kernels are \Matern{}: $\hat\phi_t(\omega) = (1 + \theta^2t^2\omega^2)^{-\gamma}$.
\item The stage transfer functions are $\hat\mu_{s,t}(\omega) =
\bigl((1 + \theta^2s^2\omega^2)/(1 + \theta^2t^2\omega^2)\bigr)^{\gamma}$ for all $s \le t$,
rational in $\omega$, with a single pole pair when $s < t$.
\end{enumerate}
The parameter $\theta$ is a change of gauge and may be normalised to $1$.
\statusT\quad\emph{(The headline of the chapter: the family of spatial statistics is, under the
hemigroup axioms, the finite-moment, rational-spectrum, exactly implementable corner of scale
space. Clause (4) is stated for integer $\gamma$ only. The draft's version asked, in a check,
which rational functions in general are admissible stage transfer functions;
\ref{prop:rational-increments} gives a class of them and the general question is open
(\ref{rem:rational-converse} shows that its first clause has no converse), while the theorem uses
only the single-pole-pair case, so the
clause is narrowed to it here rather than left with an open quantifier. \emph{The canonical gauge}
is a hypothesis since 2026-09-19 (R176): clauses (3) and (4) speak of the kernels and the stages
at the parameter $t$, which is a statement about the family's gauge as well as about its exponent,
and the declaration has read the gauge from the start. The four clauses are named
by one declaration, \texttt{SpatialLine.matern\_theorem}, which assembles the four proved
declarations over one profile and one family of kernels; it proves nothing new. Two residuals of
that assembly are recorded here rather than hidden in it. Clauses (1)--(3) are equivalent at a
\emph{given} pair $(\gamma,\theta)$ of positive reals, not merely as existential statements: the
smoothness and range read off the profile are the ones appearing in the Thorin atom and in the
transform. That is what the theorem means --- clause (4) reuses the $\theta$ of clauses
(1)--(3), which the existential reading would leave unbound --- and it costs nothing, both
directions of each equivalence producing the same pair. Clause (4) is stated at each integer
$\gamma$ and each range $\theta$, so the four are not a single chain of equivalences at one
$\gamma$; that is the theorem's own restriction of the clause to integer $\gamma$. And clause (3) is about the
endpoint kernels while clause (4) is about the increments, so the bundling reads the increment
family at its endpoint pair; the cascade law relating the two is a hypothesis of the bundle, as it
is of \texttt{matern\_rational}. Clause (4) is formalised twice. One declaration is the
equivalence with clause (3), at the closed form
$\bigl((1+\theta^2s^2\omega^2)/(1+\theta^2t^2\omega^2)\bigr)^{m}$; the other reads that closed
form as a ratio of polynomials, the denominator a power of one quadratic strictly positive on
the line --- so its two roots are a single conjugate pair off the real axis, which is the
"rational in $\omega$ with a single pole pair" of the clause's own words. The closed form on its
own does not carry it, and the integer restriction is what makes it statable: at a non-integer
exponent there is no ratio of polynomials at all. The pole-pair declaration reads the increments
at $0 < s \le t$, the collapsed left endpoint being the endpoint-kernel clause (3) rather than
this one. The closing sentence is formalised at the \emph{profile},
$k_{\gamma,\theta}(x) = k_{\gamma,1}(x/\theta)$ on $(0,\infty)$: it is a statement about the
family --- the $(\gamma,\theta)$ member is the $(\gamma,1)$ member with the displacement axis
rescaled --- and it is false without $\theta > 0$. The frequency-side identity
$F_{\gamma,\theta}(\omega) = F_{\gamma,1}(\theta\omega)$ is named beside it and is definitional,
both sides being $\gamma\log(1+\theta^2\omega^2)$; on its own it would carry none of the
sentence's content.)}
\end{theorem}

% CHANGED (proof of record 2026-09-09): rewritten to the machine-checked route, in three
% places, and one arithmetic slip is corrected.
%  (a) (1) iff (2) no longer goes through prop:thorin-subclass and Bernstein's theorem. The
%      checked proof evaluates both sides: the forward direction is the integral of
%      prop:matern-exponent(1) against a one-point measure, the backward direction is
%      uniqueness of the pair. That spends A3's uniqueness clause where the printed route
%      spent A11, and it is the cheaper of the two -- clause (2) is stated here as the
%      representation eq:thorin at a Dirac, not as membership of a named class, so the class
%      is not needed to read it.
%  (b) CORRECTION. The pair of gamma log(1 + theta^2 omega^2) was printed as
%      (0, 2 gamma theta^{-1} e^{-x/theta} x^{-1} dx). The factor theta^{-1} is spurious:
%      prop:matern-exponent(1) says the profile at range theta is k(x) = 2 gamma e^{-x/theta},
%      and the \Levy{} measure is k(x) x^{-1} dx, so the pair is
%      (0, 2 gamma e^{-x/theta} x^{-1} dx). The Lean statement carries the correct profile and
%      the error was in this proof alone; the draft's Theorem 10.3 should be corrected with it.
%  (c) THE STEP THE PRINTED PROOF SKIPPED. Uniqueness of the pair identifies the \Levy{}
%      MEASURES, hence the profiles only almost everywhere, while clauses (1) and (3) assert a
%      pointwise identity on (0,infinity) (R9's narrowing). The passage between the two is not
%      a measure-theoretic triviality and is now written out; it is ours and [T].
%  (d) (3) iff (4): the printed proof recovers the endpoint transform "by letting s -> 0",
%      which is not needed -- the increment formula is stated at s = 0 and is read there. What
%      IS needed, and was not printed, is the value at the collapsed pair s = t = 0.
\begin{proof}\leanok
\emph{(1) if and only if (2).} With $k(x) = 2\gamma e^{-x/\theta}$ and $a = 0$, take
$U = 2\gamma\,\delta_{1/\theta}$. Then \ref{eq:thorin} reads
$\tfrac12\cdot 2\gamma\log(1 + \omega^2\theta^2) = \gamma\log(1+\theta^2\omega^2)$, which is the
exponent computed in \ref{prop:matern-exponent}(1); so (1) implies (2) with that one-point
measure, the factor $2$ of the folded profile cancelling the $\tfrac12$ of \ref{eq:thorin}.
Conversely, a one-point $U = 2\gamma\,\delta_{1/\theta}$ makes the right-hand side of
\ref{eq:thorin} that same exponent, so $F(\omega) = \gamma\log(1+\theta^2\omega^2)$, and the
next paragraph's uniqueness argument returns the single exponential.

\emph{(1) implies (3)} is \ref{prop:matern-exponent}(1)--(2) after the gauge change
$t \mapsto \theta t$, which \ref{prop:matern-exponent}(1) supplies directly at range $\theta$.

\emph{(3) implies (1).} Reading the transform at the canonical scale $t = 1$ and taking
logarithms turns clause (3) into $F(\omega) = \gamma\log(1+\theta^2\omega^2)$; the other scales
say nothing further. The computation in \ref{prop:matern-exponent}(1) shows that this exponent
has the pair $(0, 2\gamma e^{-x/\theta}x^{-1}dx)$, and the pair of an exponent is unique
(\ref{prop:fourier-toolbox}(3)), so $a = 0$ and the \Levy{} measure is that one. Uniqueness of the
pair is an identity of \emph{measures}, so it gives $k(x) = 2\gamma e^{-x/\theta}$ for almost
every $x > 0$; the pointwise identity on $(0,\infty)$ follows because $k$ is nonincreasing there
while $2\gamma e^{-\cdot/\theta}$ is continuous. Indeed the set on which the two agree is
co-null, hence meets every subinterval of $(0,\infty)$; approaching a given $x > 0$ from the
left through that set gives $k(x) \le 2\gamma e^{-y/\theta}$ for points $y < x$ arbitrarily
close to $x$, whence $k(x) \le 2\gamma e^{-x/\theta}$, and approaching from the right gives the
reverse inequality. (Monotonicity of the comparison function is not used; continuity of one
side and monotonicity of the other suffice.)

\emph{(3) if and only if (4)}, for integer $\gamma$, is \ref{prop:matern-exponent}(2) read in
both directions. Forwards, the cascade law makes the increment transfer function the ratio of
the two endpoint transforms, the denominators being nonvanishing; at the collapsed pair
$s = t = 0$ the same law reads $\hat\mu_{0,1} = \hat\mu_{0,1}\hat\mu_{0,0}$ with
$\hat\mu_{0,1} > 0$, so $\hat\mu_{0,0} = 1$, which is clause (4) there. Backwards, clause (4) at
$s = 0$ is the endpoint transform of clause (3); no limit in $s$ is needed.
\end{proof}

% CHANGED (review 2026-09-09, Q11): this corollary is new; its statement and proof are the
% closing consequence of thm:matern, moved out so that that theorem's \lean tag covers all of
% what it asserts. No Lean is scheduled here, the realization clause resting on
% prop:matern-cascade (Chapter 14, no Lean); the node is therefore untagged, like the Chapter 14
% nodes it depends on, rather than carrying a tag that covers part of it.
\begin{corollary}[the \Matern{} families of integer index are the ones realized by identical first-order sections]
\label{cor:matern-realization}
\uses{thm:matern,prop:matern-exponent,prop:matern-density,prop:matern-cascade,prop:kernel-regularity}
\notes{gamma-matern-corner}
% CHANGED 2026-09-19 (R190; the author's decision D2, a widening authorized at the second external
% review round of module B): THE STATEMENT IS WIDENED, and it becomes an equivalence of three
% clauses. What it said until today prescribed the pair's transfer function,
% (1 + theta^2 s^2 omega^2)/(1 + theta^2 t^2 omega^2), so that the converse direction was clause
% (4) of thm:matern restated and the three summaries of the article -- the abstract, section 1 and
% the conclusions, all of which say only "cascades of identical first-order recursive filters, run
% forward and backward" -- claimed more than the corollary proved (referee B finding 3, referee E
% section 2.2). The pair is now arbitrary, (1 + p^2 omega^2)/(1 + q^2 omega^2), and the zero is
% FORCED to vanish at the stage from the origin: the kernel from the origin is absolutely
% continuous (prop:kernel-regularity), so its transform tends to 0 at infinity, while the ratio
% tends to (p/q)^{2 gamma}. That is the referee's five-line argument, checked here. The hypothesis
% is weakened at the same time, from every stage to ONE kernel from the origin, which is what the
% argument reads; the \uses edge to prop:kernel-regularity is added for it. No Lean is scheduled
% here and none is affected. The new clause (3) is where the widening sits; clauses (1) and (2)
% say what the node said before, with the pair no longer prescribed.
% CHANGED (2026-09-19, article mirror of module B, R174): "knot ladder" is glossed at this, its
% first use in the blueprint (the register pass's item 5; the term is chapter 14's vocabulary and
% was used here two chapters before prop:knot-exactness introduces the knots). A gloss, no change
% of content.
% CHANGED (2026-09-19, external review of module B, R170): NARROWED AGAIN, to the integer index,
% and the section is defined in the statement. R168 narrowed the class of sections and missed
% the index: a finite cascade of first-order sections has a rational transfer function, and the
% Matern stage of non-integer index is not rational, so "the Matern family" was false for
% gamma not an integer (prop:matern-cascade assumes integer gamma). The moments and the tail are
% now stated as what they are, consequences, and not as part of what is unique. Found by the
% referee. The statement read: "The Matern family is the unique admissible family with every
% moment finite, exponential tails, and an exact realization by a cascade of identical
% first-order forward--backward sections at every knot ladder."
Let $(\Phis{s}{t})$ be an admissible family in its canonical gauge and let $\gamma$ be a positive
integer. Call a symmetric kernel \emph{a cascade of $\gamma$ identical forward--backward pairs of
first-order sections} when its transform is
$\bigl((1+p^2\omega^2)/(1+q^2\omega^2)\bigr)^{\gamma}$ for some $p, q \ge 0$, the pair being the
first-order forward--backward section pair of \ref{prop:matern-cascade} with pole range $q$ and
zero range $p$. Then the following are equivalent.
\begin{enumerate}
\item The family is \Matern{} of index $\gamma$: its exponent is
$F(\omega) = \gamma\log(1+\theta^2\omega^2)$ for some $\theta > 0$.
\item Every stage is such a cascade: for all $0 \le s < t$ the increment $\mu_{s,t}$ is a cascade
of $\gamma$ identical pairs, with $p$ and $q$ depending on the pair of scales.
\item One kernel from the origin is such a cascade: for some $t > 0$ the kernel $\phi_t$ is.
\end{enumerate}
In (1) the pair of the stage from $s$ to $t$ is $(p,q) = (\theta s, \theta t)$; in (3) the zero
necessarily vanishes, $p = 0$ and $q = \theta t$. Such a family has every moment finite and
exponential tails.
% CHANGED (2026-09-18, module B register pass, R168): NARROWED. The clause read "an exact
% realization by first-order forward--backward sections at every knot ladder", which is not
% unique to the Matern family: by prop:rational-increments(1) every finite mixture of Matern
% members with integer exponent weights is realized exactly by first-order forward--backward lag
% sections at every knot ladder, and it has every moment finite and exponential tails as well.
% The proof below always argued from "a cascade of gamma such sections", all with one transfer
% function, which is clause (4) of thm:matern; the statement now says what the proof proves.
% CHANGED 2026-09-19 (R176, referee B finding 2): the canonical gauge is a hypothesis here too. The
% transfer function (1 + theta^2 s^2 omega^2)/(1 + theta^2 t^2 omega^2) prescribes how the pole and
% the zero depend on the two scales, so both directions read the gauge; in a gauge chi =/= id the
% same family's stage carries chi(s) and chi(t). Safe direction (a hypothesis added).
\statusT\quad\emph{(The engineering reading of \ref{thm:matern}, and the reason the family is
the one spatial statistics and recursive filtering both arrive at. It is stated as a corollary
rather than as a clause of the theorem because its realization half is proved in
Chapter~\ref{ch:implementation}, where the sections are constructed; no formalisation is
scheduled for it.
\textbf{The pair is not prescribed} (2026-09-19, R190, the author's decision D2). Until today the
statement fixed the pair's transfer function to $(1+\theta^2s^2\omega^2)/(1+\theta^2t^2\omega^2)$,
which ties the pole and the zero to the two scales in advance; with that prescription the converse
is clause (4) of \ref{thm:matern} restated. The statement now lets the pair be any
$(1+p^2\omega^2)/(1+q^2\omega^2)$, so that what is characterised is the \emph{shape} of the
realization --- a cascade of identical first-order forward--backward pairs --- and not a formula
already containing the answer. The three summaries of the article that announce the corollary are
then correct as they stand.
\textbf{Where the hypothesis bites, and how little of it is needed.} Clause (3) asks for
\emph{one} kernel from the origin, so the stages between two positive scales may be dropped from
the hypothesis altogether; and the number of pairs cannot vary from stage to stage, by the proof's
last paragraph. What may not be dropped is the word \emph{identical}: an admissible family whose
kernel from the origin is a cascade of first-order forward--backward pairs of different ranges
need not be \Matern{}, and the second example of \ref{rem:rational-converse} is one ---
$F(\omega) = 2\log(1+\omega^2) - \log(1+\omega^2/4)$, whose kernel from the origin at scale $t$
has transform
$\bigl[(1 + t^2\omega^2/4)/(1 + t^2\omega^2)\bigr]\cdot(1 + t^2\omega^2)^{-1}$, a cascade of two
pairs, one carrying a zero of range $t/2$ and one a pure pole of range $t$. The example the
external review gave is the same phenomenon with three pairs: $k(x) = 6e^{-x} - 2e^{-2x}$, which is
positive ($6e^{-x} \ge 2e^{-x} > 2e^{-2x}$) and nonincreasing
($k'(x) = -2e^{-x}(3 - 2e^{-x}) < 0$) and satisfies both integrability conditions, has
$F(\omega) = 3\log(1+\omega^2) - \log(1+\omega^2/4)$ by \ref{prop:matern-exponent}(1) at the two
rates. Neither family has a single-exponential profile.
\textbf{Not claimed}: that the pair is realized by a \emph{stable} recursion for every $(p,q)$ ---
the realization of \ref{prop:matern-cascade} is read at $p < q$, which the proof derives, and the
sections there are the ones with positive coefficients; and anything about non-integer index,
where the stage is not rational at all.)}
\end{corollary}

\begin{proof}
\emph{(1) implies (2).} A \Matern{} family of positive integer index $\gamma$ and range $\theta$
has, by \ref{prop:matern-cascade} read in the gauge $t \mapsto \theta t$, the stage transfer
functions $\bigl((1+\theta^2s^2\omega^2)/(1+\theta^2t^2\omega^2)\bigr)^\gamma$ at every pair of
scales, realized by $\gamma$ identical forward--backward pairs of first-order sections; so every
stage is a cascade of the kind named, with $(p,q) = (\theta s, \theta t)$, and $p < q$.

\emph{(2) implies (3)} is the stage from $s = 0$, which in the canonical gauge is the kernel
$\phi_t = \mu_{0,t}$ from the origin, at any $t > 0$.

\emph{(3) implies (1).} Suppose $\hat\phi_t(\omega) = \bigl((1+p^2\omega^2)/(1+q^2\omega^2)\bigr)
^{\gamma}$ for some $t > 0$ and some $p, q \ge 0$. First $q > p$. The transform of a probability
measure is bounded by $1$ in modulus, so the ratio $(1+p^2\omega^2)/(1+q^2\omega^2)$ is bounded on
the line, which forces $q \ge p$; and $q = p$ would make $\hat\phi_t \equiv 1$, hence
$\phi_t = \delta_0$ and $\Phis{0}{t} = \mathrm{Id}$, against (ND). Next $p = 0$. By
\ref{prop:kernel-regularity} the kernel $\mu_{0,t}$ from the origin is absolutely continuous, so
its transform tends to $0$ at infinity --- the Fourier transform of an integrable function does,
and $\hat\phi_t$ is that transform --- while
\[
  \Bigl(\frac{1 + p^2\omega^2}{1 + q^2\omega^2}\Bigr)^{\gamma}
  \ \longrightarrow\ \Bigl(\frac pq\Bigr)^{2\gamma}
  \qquad (|\omega| \to \infty),
\]
a positive number unless $p = 0$. So $p = 0$ and $q > 0$. Reading the canonical gauge,
$e^{-F(t\omega)} = (1+q^2\omega^2)^{-\gamma}$ for every $\omega$, that is
$F(t\omega) = \gamma\log(1+q^2\omega^2)$; substituting $\omega \mapsto \omega/t$ gives
$F(\omega) = \gamma\log(1+\theta^2\omega^2)$ with $\theta := q/t > 0$, which is clause (1).

\emph{The index is not free.} Suppose the kernels from the origin at two scales are cascades of
$\gamma$ and $\gamma'$ identical pairs. Each gives, by the paragraph just proved, an exponent
$\gamma\log(1+\theta^2\omega^2) = \gamma'\log(1+\theta'^2\omega^2)$; dividing by $\log|\omega|$
and letting $|\omega| \to \infty$ gives $2\gamma = 2\gamma'$. So the number of pairs is the same
at every scale, and it is the index.

\emph{The consequences.} Every moment of such a family is finite by \ref{prop:matern-exponent}(3),
and the tail is exponential by \ref{prop:matern-density}.
\end{proof}

% draft: Proposition 10.4
% twin: hcs prop:thorin-subclass (ledger A21 there); clauses (4) and (5) are new.
% SKELETON NOTE (review 2026-09-09, R19): CLAUSE (3) IS UNTYPED, and deliberately. The phase B
% declaration Skeleton.thorin_subclass_ggc carried the class of symmetric extended generalized
% gamma convolutions as a parameter whose defining hypothesis was clause (1)'s own form, so it
% was provable in [T] and carried none of A17 -- an [A] clause discharged by a tautology.
% Defining the class instead from A17's own statement-as-used is not available either: that
% statement is a MOMENT GENERATING function on real sigma, and the laws this node is about have
% none (the Student-t law has no exponential moment at any sigma =/= 0), while the reading that
% does apply to them is the same formula at sigma = i omega, whose +-theta pairing the ledger
% explicitly does NOT carry. So the declaration was withdrawn; the \lean tag below covers
% clauses (1)-(2), (4) and (5). Nothing downstream consumes clause (3) by name.
% CHANGED (proof 2026-09-10, wave 6): CLAUSE (3) IS MOVED OUT into lem:thorin-ggc below and
% this node is leanok. The reason is R27's, at a statement rather than at a declaration, and it
% is the reason the wave 5 merge gave for lem:polya-student-exclusion: all three declarations
% this tag names are now proved, and the one clause with no Lean -- clause (3), untyped since
% the review for the reasons in the SKELETON NOTE above -- would have held the whole node
% notready by itself. Split, each part is honest: what this node claims is what its three
% declarations prove, and clause (3) is an [A] lemma of its own on ledger A17, with no Lean
% scheduled and the reason written into it. THE CLAUSE NUMBERING IS UNCHANGED -- (1), (2), (4),
% (5), with (3) absent and named where it went -- because renumbering would silently invalidate
% every existing reference to prop:thorin-subclass(4) and (5), of which prop:student-t's
% annotation alone has two. The split adds one uses edge, from the new lemma to this node, and
% none here.
% CHANGED (2026-09-11, R79, under the safe-direction delegation): clause (5) closed with "with the
% Matern, Student-t and stable families in the smallest class and the Cin rays outside the middle
% one" -- a membership sentence of exactly the shape R44 found at prop:bridge-strictness(4), at a
% \leanok node whose three declarations prove the representation equivalence, the bridge identity
% and the strict inclusion, and nothing about any named family. The Matern part is this node's own
% one-atom case, the stable part is prop:stable-family's Thorin density, the Cin part is
% prop:bridge-strictness(4); the Student-t part is prop:student-t(3), which is sorry on A18. The
% sentence is moved to the paragraph after the proof with each membership sent to the node that
% states it, and the clause keeps the chain of strict inclusions, which is what thorin_strictness
% proves. No proof changed.
% CHANGED (2026-09-11, R120, under the safe-direction delegation): the tag now covers clause (4)'s
% second sentence. thorin_bridge proved only the forward map, while the clause also says that
% U_I |-> U is a linear bijection of the two domains and that the bridge restricts to a bijection
% of the causal Thorin subclass onto the symmetric one. Three declarations are added:
% thorin_bridge_onto (every symmetric Thorin member is the bridge image of a causal one, with
% b_0 = 2a and U twice the image of U_I), thorin_bridge_measure_injective and
% thorin_bridge_measure_linear; all three print Lean core. The printed "hence every symmetric
% Thorin measure is twice the image of a causal one" rested on the term-by-term correspondence of
% the integrability conditions; the proof of record is rewritten to the checked route, which
% CONSTRUCTS the causal datum and COMPUTES its exponent (the causal Frullani integral under U_I),
% and bounds the causal integrability fields by the spatial Thorin integral at omega = 1. The
% term-by-term correspondence is kept as a marked remark. No statement changed.
% CHANGED (2026-09-19, article mirror of module B, R174): the paragraph between clause (2) and
% clause (4) no longer carries ledger ids or formalisation scheduling, which is the register pass's
% item 11: a statement of record says what is true, and where an entry binds and what is scheduled
% belong to the annotation. It still says that the third member of the equivalence is
% lem:thorin-ggc, why it is stated apart (the reading of Thorin's definition, not Bernstein's
% theorem) and that the numbering skips (3) -- a reader needs the last of these to read the
% numbering, and it is a fact and not an apology. The entry (A17), the contrast with A11 and the
% scheduling are now the first sentence of the annotation. No clause and no Lean changes; the
% numbering itself is unchanged, for the reason the wave 6 block above gives.
% CHANGED 2026-09-19 (R175, second external review round of module B, referee B finding 1): THE
% INTEGRABILITY CONDITION OF eq:thorin IS RESTATED, and the description of the bridge's target in
% clause (4) -- now clause (3) -- is NARROWED. The condition printed until today,
% int_(0,1] log(1/theta) U + int_(1,infty) theta^{-2} U < infinity, does not control U near
% theta = 1 (log(1/theta) vanishes there and the second integral starts above it): the referee's
% U = sum_n delta_{1 - 1/n^2} satisfies it with sum log(1/theta_n) = log 2, is not locally finite,
% and makes the right-hand side of eq:thorin infinite at every omega =/= 0. Clause (2) was not
% false -- the identity it asserts forces the strong condition anyway -- but three things around it
% were: (a) the (2) implies (1) step "which the same three regimes make finite", which needs
% U((0,1]) < infinity and not the printed pair of integrals; (b) "the two domains of Thorin
% measures" in clause (4), the printed domain being strictly larger than the image of the causal
% one; (c) the term-by-term correspondences in the remark after the proof and in the proof of
% lem:thorin-ggc, Bondesson's (7.1.2) being int (1+theta^2)^{-1} V + int_{|theta|<=1} |log theta^2| V,
% which DOES control V on compacta. The condition is now the single integral
% int log(1 + theta^{-2}) U < infinity, which is what every proof of the node uses, what the Lean
% reads off one frequency (thorin_lintegral_ne_top, then thorin_measure_facts for the two
% integrals and for U((0,1]) < infinity), and what corresponds to the causal condition and to
% Bondesson's; it is equivalent to the printed pair together with local finiteness on (0,infty).
% Clause (3) (formerly (4)) no longer says "the two domains": it says linear, injective, and onto
% the Thorin measures of the symmetric Thorin members, which is exactly what the three
% declarations thorin_bridge_measure_linear, thorin_bridge_measure_injective and
% thorin_bridge_onto prove (the last is stated for a U representing the exponent of an SDProfile,
% not for an arbitrary U in a domain). Nothing in the Lean is false or vacuous under the weak
% condition: no declaration takes it as a hypothesis and concludes from it --
% thorin_subclass_representation asserts it in the forward direction and carries it inertly in the
% converse, where the identity does the work.
% CHANGED 2026-09-19 (R175, referee B finding 11): THE CLAUSES ARE RENUMBERED 1, 2, 3, 4. The gap
% at (3) dates from the wave 6 split (the block above), which kept the numbering so as not to
% invalidate references, and the author's standing decision was to renumber at the draft sync. The
% draft is frozen (ADR-0008), so the sync will not come; every reference to the old (4) and (5) in
% the blueprint and in paper-b is updated in the same commit.
\begin{proposition}[the symmetric Thorin subclass]\label{prop:thorin-subclass}
\uses{lem:selfdecomposable-exponents,prop:matern-exponent,lem:bridge-exponents,prop:bridge-families,prop:bridge-strictness,prop:fourier-toolbox,prop:laplace-uniqueness-locally-finite,lem:quadratic-growth}
\lean{SpatialLine.thorin_subclass_representation, SpatialLine.thorin_bridge,
SpatialLine.thorin_bridge_onto, SpatialLine.thorin_bridge_measure_injective,
SpatialLine.thorin_bridge_measure_linear, SpatialLine.thorin_strictness}\leanok
\notes{gamma-matern-corner}
For an admissible $F$ with data $(a,k)$ the following are equivalent.
\begin{enumerate}
\item $k$ is completely monotone.
\item $F$ has a symmetric Thorin representation
\begin{equation}\label{eq:thorin}
  F(\omega) = a\,\omega^2
   + \tfrac12\int_{(0,\infty)} \log\Bigl(1 + \frac{\omega^2}{\theta^2}\Bigr)\,U(d\theta),
  \qquad U \ge 0, \quad
  \int_{(0,\infty)}\log\Bigl(1 + \frac{1}{\theta^2}\Bigr)\,U(d\theta) < \infty .
\end{equation}
\end{enumerate}
A third member of the equivalence --- that the kernels are symmetric extended generalized gamma
convolutions in Thorin's sense --- is stated apart, as \ref{lem:thorin-ggc}, because it rests on
the reading of Thorin's definition rather than on Bernstein's theorem.

Moreover:
\begin{enumerate}
\item[(3)] \emph{(The bridge on Thorin subclasses.)} If $F_{\mathrm I}$ is causally admissible
with completely monotone profile $k_{\mathrm I} = \int e^{-\theta u}\,U_{\mathrm I}(d\theta)$,
then $F_{\mathrm I}(\cdot^2/2)$ is a symmetric Thorin member with $a = \tfrac12 b_0$ and $U$
equal to \emph{twice} the image of $U_{\mathrm I}$ under $\theta \mapsto \sqrt{2\theta}$; the
map $U_{\mathrm I} \mapsto U$ is linear and injective, and carries the causal Thorin measures
onto the Thorin measures of the symmetric Thorin members, so the subordination bridge restricts
to a linear bijection of the causal Thorin subclass onto the symmetric one.
\item[(4)] \emph{(Strictness the other way.)} The subordinated slice $\mathcal{S}$ of
\ref{prop:bridge-strictness} is strictly larger than the symmetric Thorin subclass: the image
of the causal ray with profile $k_{\mathrm I} = \mathbf 1_{(0,\tau)}$ has the explicit profile
$k(x) = 2\,\erfc\bigl(x/\sqrt{2\tau}\bigr)$, which is not completely monotone. So the symmetric
Thorin subclass is strictly inside $\mathcal{S}$, which is strictly inside the admissible cone.
\end{enumerate}
\statusA\quad\ledger{A11}\quad\emph{(\textbf{Assignment.} The third member of the equivalence is
\ref{lem:thorin-ggc}: it is the one that rests on \ledger{A17} rather
than on \ledger{A11}, and the one for which no formalisation is scheduled; the statement above
says where it went and says no more, the entry and the scheduling being recorded here (2026-09-19,
R174). \ledger{A11} carries
Bernstein's theorem, which is clause (1) if and only if (2) --- a completely monotone $k$ is
$\int e^{-\theta x}U(d\theta)$ for a unique positive $U$. What it does not carry: the computation
$\int_0^\infty (1-\cos\omega x)e^{-\theta x}x^{-1}dx = \tfrac12\log(1+\omega^2/\theta^2)$ and
the Tonelli that turns it into \ref{eq:thorin}; the integrability correspondence; and clauses
(3) and (4) entirely, which are \textbf{[T]} computations on \ref{eq:bridge-profile}. An atom
of $U$ at $\theta = 0$ is excluded because it would make $k$ bounded below and
$\int_1^\infty k(x)x^{-1}dx$ infinite; that too is \textbf{[T]}. The formalisation admits
\ledger{A11} as the \emph{existence equivalence only}: the uniqueness of the Thorin measure is
\ref{prop:laplace-uniqueness-locally-finite}, which this article proves, so the uniqueness half of
the entry is never cited. The declarations for (1) if and only if (2), for the forward map of
(3) and for (4) spend that one name and \ledger{A3}'s uniqueness clause, and nothing else:
clauses (3) and (4) reach \ledger{A11} at its two opposite directions --- (3) at the easy one, a
Laplace transform is completely monotone, and (4) at the forward one, in the contradiction ---
and reach \ledger{A3} through the identification of a quantified profile with the subordinated
one, which is an almost-everywhere identity upgraded to a pointwise one by the continuity of the
transform. The three declarations for the rest of (3) --- that every symmetric Thorin member is
the image of a causal one, and that $U_{\mathrm I} \mapsto U$ is linear and injective --- spend
no interface: the causal datum is constructed and its exponent computed. \emph{The integrability
condition of \ref{eq:thorin} and the formal statement.} \texttt{thorin\_subclass\_representation}
asserts, beside the identity, the pair of integrals
$\int_{(0,1]}\log(1/\theta)\,U(d\theta) + \int_{(1,\infty)}\theta^{-2}U(d\theta) < \infty$, which
is what the printed condition was until 2026-09-19 (R175). The displayed condition is the
equivalent one the proofs use: it is the Thorin integral at $\omega = 1$, which
\texttt{thorin\_lintegral\_ne\_top} reads off the finiteness of the exponent, and
\texttt{thorin\_measure\_facts} derives from it the two integrals above \emph{and} the
$\sigma$-finiteness and local finiteness ($U((0,1]) < \infty$) that the converse direction
needs. The pair of integrals alone does not give the last of these, which is why it is no longer
what the display prints. \emph{The formal statement of (2) says one
thing more than the display:} \texttt{thorin\_subclass\_representation} bundles into the same
equivalence that the folded measure $U$ of \ref{eq:thorin} is \emph{unique}, so the Lean is
strictly stronger than the printed clause, which asserts existence only. The uniqueness is
\ref{prop:laplace-uniqueness-locally-finite} read at the tail of $k$, proved here and free of
\ledger{A11}'s own uniqueness half; the display does not carry it because no consumer in this
article cites it, and the paragraph above says why the entry's uniqueness clause is never
spent.
\textbf{Attribution of clause (3)} (2026-09-19, R196, left by batch S2; the page was read at the
held copy this session). Both directions, with the correspondence of the measures, are
Bondesson's Theorem~7.3.1, p.~115: if $\varphi$ is the moment generating function of a
generalized gamma convolution then $\psi(s) = \varphi(s^2/2)$ is that of a symmetric extended
generalized gamma convolution, and conversely; and the proof on pp.~115--116 computes the
spectral measure of $\psi$ as the image of that of $\varphi$ under $t \mapsto \pm\sqrt{2t}$, the
left extremity $a$ of $\varphi$ becoming the coefficient $a/2$ of $s^2$ in $\psi$. Read at
$s = i\omega$ --- which is \ref{lem:thorin-ggc} --- that is this clause. What belongs to this
article is the folded normalisation, $U$ twice the restriction of the spectral measure to
$(0,\infty)$; the linearity and injectivity of the measure map; the verification that the
constructed causal datum is causally admissible; and that all of it is machine-checked.)}
\end{proposition}

% CHANGED (2026-09-11, draft sync, R103): the wave 3 block below stood at prop:moments-tails,
% whose proof it does not describe; it is this node's (1) <=> (2) and the uniqueness of U, and is
% moved here verbatim but for "at theta = 1", which reads "at x = 1" (R106, next comment).
% CHANGED (proof of record 2026-09-09, wave 3): the proof of the equivalence of clauses (1)
% and (2) is rewritten to the machine-checked route (repo CLAUDE.md, author's decision of the
% same day). Three steps the printed argument left implicit are now stated, because the
% formalisation had to supply each of them: sigma-finiteness of U, which the transform's
% finiteness at x = 1 gives; that the integrability condition of eq:thorin is a
% CONSEQUENCE of the identity and not a further hypothesis; and, in the converse direction,
% that the reconstructed profile k' satisfies the \Levy{} condition at all, which follows by
% averaging the frequency identity against a standard Gaussian law. Uniqueness of U is now a
% paragraph of its own and runs on prop:laplace-uniqueness-locally-finite rather than on the
% uniqueness half of A11, which is therefore NOT admitted as an axiom; the annotation records
% that. Clauses (4) and (5) are untouched and remain unformalised.
% CHANGED (2026-09-11, draft sync, R106): (1) implies (2) said U is sigma-finite "its transform
% being finite at theta = 1". The transform int e^{-theta x} U(d theta) is a function of x (theta
% is the integration variable), and what is finite is its value k(1); so "at x = 1", as the
% draft already read.
% CHANGED (proof of record 2026-09-10, wave 6): clauses (4) and (5) are rewritten to the
% machine-checked route (repo CLAUDE.md, author's decision of 2026-09-09).
% (4) gains two paragraphs, both of them steps the printed argument left implicit and the
% formalisation had to supply. One is sigma-finiteness of U_I, which the Laplace hypothesis at
% u = 1 gives for free and which the Tonelli needs. The other is the IDENTIFICATION of k: the
% exponent determines the profile only up to a null set, while complete monotonicity is a
% statement about every derivative at every point, so the passage from almost everywhere to
% everywhere has to be made, and it is thm:matern's -- one side nonincreasing, the other a
% continuous Laplace transform. Neither step changes the route; both were invisible in the
% printed proof and are the reason wave 5 re-priced the clause.
% (5)'s non-complete-monotonicity half is rewritten ENTIRELY. The printed argument compares an
% exponential floor against the ASYMPTOTIC erfc(y) ~ sqrt(2 tau/pi) y^{-1}e^{-y^2}, which is
% Watson's lemma and is more than the conclusion needs. The checked route is a bound at one
% point: the floor c e^{-nx} comes from a single window of Bernstein's measure carrying positive
% mass, the ceiling from the elementary comparison e^{-v^2} <= e^{-yv} on (y, infinity), and the
% two are contradicted at an explicit x. The printed route is kept as a marked remark.
\begin{proof}[Proof of the parts held \textbf{[T]}]
\leanok
\emph{(1) implies (2).} By A11, $k(x) = \int_{[0,\infty)} e^{-\theta x}\,U(d\theta)$ for a
positive measure $U$ on $[0,\infty)$; an atom at $\theta = 0$ would make $k$ bounded below by a
positive constant, so that $\int_1^\infty k(x)x^{-1}dx = \infty$, against admissibility. So $U$
lives on $(0,\infty)$, and it is $\sigma$-finite, its transform being finite at $x = 1$.
Tonelli with the computation of \ref{prop:matern-exponent}(1) gives
\[
  \int_0^\infty (1-\cos\omega x)\,k(x)\,\frac{dx}{x}
  = \int_{(0,\infty)}\Bigl(\int_0^\infty (1-\cos\omega x)\,e^{-\theta x}\,\frac{dx}{x}\Bigr) U(d\theta)
  = \tfrac12\int_{(0,\infty)}\log\Bigl(1 + \frac{\omega^2}{\theta^2}\Bigr)U(d\theta),
\]
which is \ref{eq:thorin}. The integrability condition of \ref{eq:thorin} is not a further
hypothesis but a consequence of that identity: read at $\omega = 1$ the identity says that the
integral is $2(F(1) - a)$, and $F(1)$ is finite for every admissible $F$ by
\ref{lem:quadratic-growth}. Three regimes read off from it what the classical form of the
condition states separately: $\log(1+\theta^{-2}) \ge 2\log(1/\theta)$ near $0$,
$\ge \tfrac12\theta^{-2}$ at infinity, and $\ge \log 2$ on $(0,1]$, so that
$\int_{(0,1]}\log(1/\theta)\,U(d\theta)$ and $\int_{(1,\infty)}\theta^{-2}\,U(d\theta)$ are
finite and $U$ is finite on $(0,1]$, hence on every compact subset of $(0,\infty)$.

\emph{(2) implies (1).} Given \ref{eq:thorin} with $U$ carried by $(0,\infty)$, set
$k'(x) := \int_{(0,\infty)} e^{-\theta x}\,U(d\theta)$, which is finite for every $x > 0$: on
$(0,1]$ the integrand is at most $1$ and $U$ is finite there, and on $(1,\infty)$ one has
$e^{-\theta x} \le C_x\theta^{-2}$ against the second integral, both facts being the three
regimes of the previous paragraph. It is the finiteness of $U$ on $(0,1]$ that is doing the work
here, and it is the reason the displayed condition is the one printed: the two integrals alone
leave $U$ free to put infinite mass just below $\theta = 1$. The displayed Tonelli, read
backwards, says that $k$ and $k'$ have the same
compensated integral $\int_0^\infty(1-\cos\omega x)\,\cdot\,dx/x$ at every frequency. That $k'$
is the profile of a symmetric \Levy{} pair at all --- that is, that
$\int_0^\infty (1\wedge x^2)\,k'(x)\,dx/x < \infty$ --- follows by averaging that identity over
$\omega$ against a standard Gaussian law: Tonelli replaces the weight $1 - \cos\omega x$ by
$1 - e^{-x^2/2}$, which lies between $\tfrac13(1 \wedge x^2)$ and $1 \wedge x^2$, so the \Levy{}
condition transfers from $k$ to $k'$. Uniqueness of the pair (\ref{prop:fourier-toolbox}(3)) then
identifies the two \Levy{} measures, hence $k$ and $k'$ almost everywhere on $(0,\infty)$, and hence
everywhere there, $k$ being nonincreasing and $k'$ continuous. So $k$ is a Laplace transform, and
A11 read the other way returns complete monotonicity.

\emph{Uniqueness of $U$.} If $U'$ on $(0,\infty)$ also satisfies \ref{eq:thorin}, the previous
paragraph applied to $U'$ gives $k(x) = \int e^{-\theta x}\,U'(d\theta)$ as well, so $U$ and $U'$
have the same Laplace transform on $(0,\infty)$ and it is finite there; by
\ref{prop:laplace-uniqueness-locally-finite}, $U' = U$. This does not spend the uniqueness half
of A11: the Laplace uniqueness the argument needs is proved in this article.

\emph{(3).} With $k_{\mathrm I}(u) = \int e^{-\theta u}\,U_{\mathrm I}(d\theta)$,
\ref{eq:bridge-profile} and Tonelli give
\[
  k(x) = \int U_{\mathrm I}(d\theta)\;2x\int_0^\infty g_u(x)\,e^{-\theta u}\,\frac{du}{u}
       = 2\int e^{-x\sqrt{2\theta}}\,U_{\mathrm I}(d\theta),
\]
by the first-passage identity used in \ref{prop:bridge-families}(2); substituting
$\theta' = \sqrt{2\theta}$ this is $\int e^{-\theta'x}\,U(d\theta')$ with $U$ twice the image
of $U_{\mathrm I}$, so $k$ is completely monotone, and the Gaussian coefficient is
$\tfrac12 b_0$ by \ref{lem:bridge-exponents}. That $U_{\mathrm I}$ is $\sigma$-finite, which the
Tonelli needs, is not a further hypothesis: $k_{\mathrm I}(1)$ is a real number, so the
transform of $U_{\mathrm I}$ is finite at $u = 1$ against a weight vanishing nowhere.

The identification of $k$ deserves a sentence, because the exponent determines it only up to a
null set. What \ref{prop:fourier-toolbox}(3) returns, applied to the subordinated exponent and to
any admissible $F$ with that exponent, is equality of the two \Levy{} measures, hence of the two
profiles almost everywhere on $(0,\infty)$; the conclusions of this clause --- complete
monotonicity, which is a statement about every derivative at every point --- are pointwise. The
passage is the one \ref{thm:matern} makes: the two profiles are nonincreasing and the Laplace
transform $\int e^{-\theta'x}U(d\theta')$ is continuous on $(0,\infty)$ by dominated convergence,
and a nonincreasing function almost everywhere equal to a continuous one equals it everywhere,
the good set meeting every subinterval on each side of a given point.

The factor $2$ is the folding convention and is
checked on the \Matern{} member: the causal Gamma data are $U_{\mathrm I} = \gamma\delta_1$,
whence $U = 2\gamma\delta_{\sqrt2}$ and \ref{eq:thorin} returns
$\tfrac12\cdot 2\gamma\log(1+\omega^2/2) = \gamma\log(1+\omega^2/2)$, which is
\ref{prop:bridge-families}(2); the image without the factor would return half of it.

\emph{The bridge is onto the symmetric Thorin subclass.} Let $F$ be admissible with data $(a,k)$
and Thorin measure $U$ on $(0,\infty)$ as in \ref{eq:thorin}. Put $b_0 := 2a$, let
$U_{\mathrm I}$ be half the image of $U$ under $\theta' \mapsto \theta'^2/2$, and let
$k_{\mathrm I}(u) := \int e^{-\theta u}\,U_{\mathrm I}(d\theta)$ for $u > 0$. The transform is
finite: the exponents satisfy
$u\theta'^2/2 - (u\theta'/2 - u/2) = \tfrac u2\bigl((\theta' - \tfrac12)^2 + \tfrac34\bigr) \ge 0$,
so $k_{\mathrm I}(u) \le \tfrac12 e^{u/2}\int e^{-u\theta'/2}\,U(d\theta')$, which the three
regimes of (1) implies (2) make finite. Hence $k_{\mathrm I}$ is nonnegative and nonincreasing on
$(0,\infty)$, and $U_{\mathrm I}$ is $\sigma$-finite. The exponent of $(b_0, k_{\mathrm I})$ is
computed rather than identified. For $u > 0$ and $\sigma \ge 0$ one has
$(1 - e^{-\sigma u})/u = \int_0^\sigma e^{-su}\,ds$, and integrating out $u$ first,
\[
  \int_0^\infty (1 - e^{-\sigma u})\,k_{\mathrm I}(u)\,\frac{du}{u}
  = \int U_{\mathrm I}(d\theta)\int_0^\sigma \frac{ds}{s + \theta}
  = \int \log\Bigl(1 + \frac{\sigma}{\theta}\Bigr)\,U_{\mathrm I}(d\theta)
\]
by Tonelli, twice. At $\sigma = \omega^2/2$ the substitution $\theta = \theta'^2/2$ turns
$\log(1 + \sigma/\theta)$ into $\log(1 + \omega^2/\theta'^2)$ and $U_{\mathrm I}$ into
$\tfrac12 U$, and $b_0\sigma = a\omega^2$, so $F_{\mathrm I}(\omega^2/2)$ is the right-hand side
of \ref{eq:thorin}, which is $F(\omega)$. At $\sigma = 1$ the display gives admissibility. The two
integrability conditions of the causal cone together say
$\int_0^\infty \min(1,u)\,k_{\mathrm I}(u)\,du/u < \infty$, and $\min(1,u) \le e\,(1 - e^{-u})$
for $u \ge 0$: for $u \le 1$ because $e\,(1 - e^{-u}) \ge e^u(1 - e^{-u}) = e^u - 1 \ge u$, and
for $u \ge 1$ because $e\,(1 - e^{-u}) \ge e - 1 \ge 1$. So the conditions are bounded by
$e\int \log(1 + 1/\theta)\,U_{\mathrm I}(d\theta)
= \tfrac e2\int \log(1 + 2/\theta'^2)\,U(d\theta')$, and $\log(1 + 2x) \le 2\log(1 + x)$, which
is $1 + 2x \le (1 + x)^2$, bounds that by $e$ times the Thorin integral of \ref{eq:thorin} at
$\omega = 1$, finite because $F(1)$ is. Finally $\sqrt{2\cdot\theta'^2/2} = \theta'$ on
$(0,\infty)$, so $U$ is twice the image of $U_{\mathrm I}$ under $\theta \mapsto \sqrt{2\theta}$.
So every symmetric Thorin member is the bridge image of a causal Thorin member, with Gaussian
coefficient and Thorin measures related as in (4); no uniqueness theorem and no part of A11 is
used.

\emph{The measure map is a linear bijection.} It is linear because taking an image measure and
multiplying by $2$ both are. It is injective: if twice the images of $U_{\mathrm I}$ and
$U_{\mathrm I}'$ under $\theta \mapsto \sqrt{2\theta}$ agree, so do the images, and the image of
either under $\theta' \mapsto \theta'^2/2$, which inverts $\theta \mapsto \sqrt{2\theta}$ on
$(0,\infty)$, returns the measure itself. Surjectivity onto the symmetric Thorin measures is the
previous paragraph.

\emph{Remark, not machine-checked.} The two conditions also correspond directly, without passing
through an exponent. Write
$C_{\mathrm I} := \int\log(1 + \theta^{-1})\,U_{\mathrm I}(d\theta)$ for the causal condition and
$C := \int\log(1 + \theta'^{-2})\,U(d\theta')$ for \ref{eq:thorin}'s. Since $U$ is twice the image
of $U_{\mathrm I}$ under $\theta \mapsto \sqrt{2\theta}$, the substitution $\theta'^2 = 2\theta$
gives $C = 2\int\log(1 + v)\,U_{\mathrm I}(d\theta)$ with $v := 1/2\theta$, while
$C_{\mathrm I} = \int\log(1 + 2v)\,U_{\mathrm I}(d\theta)$; and
$\log(1+v) \le \log(1+2v) \le 2\log(1+v)$, the second inequality being $1 + 2v \le (1+v)^2$. So
$C_{\mathrm I} \le C \le 2C_{\mathrm I}$, and each condition holds exactly when the other does.
The checked route does not compare the conditions at all: it reads both off the finiteness of one
exponent value.

\emph{(4).} The causal ray has $k_{\mathrm I} = \mathbf 1_{(0,\tau)}$, and
\ref{eq:bridge-profile} gives $k(x) = 2x\int_0^\tau g_u(x)\,du/u$. Substituting $v = x^2/2u$,
so that $du/u = -dv/v$ and $g_u(x) = \sqrt{v/\pi}\,x^{-1}e^{-v}$,
\[
  \int_0^\tau g_u(x)\,\frac{du}{u}
  = \frac{1}{\sqrt\pi\,x}\int_{x^2/2\tau}^\infty v^{-1/2}e^{-v}\,dv
  = \frac{1}{x}\,\erfc\Bigl(\frac{x}{\sqrt{2\tau}}\Bigr),
\]
so $k(x) = 2\,\erfc(x/\sqrt{2\tau})$, identified pointwise by the passage of clause (3) ---
$\erfc$ is continuous, being an integral over a moving half-line of an integrable function.

It is not completely monotone, and one bound settles it. Suppose it were; then by A11 it is
$\int e^{-\theta x}\,V(d\theta)$ for a positive $V$ on $[0,\infty)$, and $V \ne 0$ because
$\erfc > 0$. Some window carries mass: $V(-\infty, n] > 0$ for an integer $n$, and that mass is
finite because the transform is. On that window $e^{-\theta x} \ge e^{-nx}$, so
\[
  k(x) \;\ge\; c\,e^{-nx}, \qquad c := V(-\infty,n] > 0,
\]
for every $x > 0$. Against that, for $y \ge 1$ the comparison $e^{-v^2} \le e^{-yv}$ on
$(y,\infty)$ gives $\int_y^\infty e^{-v^2}dv \le e^{-y^2}/y \le e^{-y^2}$, so
$\erfc y \le 2e^{-y^2}$ and $k(x) \le 4e^{-x^2/2\tau}$ once $x \ge \sqrt{2\tau}$. Take
$x \ge \max\bigl(\sqrt{2\tau},\,2\tau(n+1),\,4/c + 1\bigr)$. The middle bound gives
$x^2/2\tau \ge (n+1)x$ and the last gives $4e^{-x} \le 4/x < c$, whence
\[
  k(x) \;\le\; 4e^{-x^2/2\tau} \;\le\; 4e^{-x}\,e^{-nx} \;<\; c\,e^{-nx} \;\le\; k(x),
\]
which is absurd. Its image is admissible by
\ref{lem:bridge-exponents}, so it lies in $\mathcal{S}$ and not in the Thorin subclass; the
other strict inclusion is \ref{prop:bridge-strictness}(4).

\emph{Remark, not machine-checked.} A second route reads the same conclusion off the
asymptotic $\erfc(x/\sqrt{2\tau}) \sim \sqrt{2\tau/\pi}\,x^{-1}e^{-x^2/2\tau}$, which says more
than the contradiction needs: an upper bound of Gaussian type at one sufficiently large $x$ is
enough, and the bound above costs a monotone comparison where the asymptotic costs Watson's
lemma.
\end{proof}

Where the named families sit in the chain of (5) is said by their own nodes, not by this one: the
\Matern{} family is the one-atom case of (2), so it lies in the symmetric Thorin subclass; the stable
family's Thorin density is \ref{prop:stable-family}; the Student-t family's membership is
\ref{prop:student-t}(3), machine-checked under the hypothesis that the causal delay law is a
generalized gamma convolution, which is the cited fact \ref{rem:student-ggc} (\ledger{A18},
unadmitted); and the
$\Cin$ rays lie outside $\mathcal{S}$ by \ref{prop:bridge-strictness}(4). (Until 2026-09-11
this sentence stood inside the strictness clause; fidelity ledger row R79. Its description of the Student-t
membership was left as ``rests on \ledger{A18} and is not yet machine-checked'' when R160 made
that clause conditional on 2026-09-15, and is corrected here; 2026-09-19, R174.)

% draft: Proposition 10.4, clause (3)
% SKELETON NOTE (2026-09-10, wave 6): NO LEAN IS SCHEDULED, and the reason is the one the
% review recorded at R19 and the wave 6 split now states at a node of its own rather than in a
% comment on somebody else's. A17's statement-as-used is a MOMENT GENERATING function on real
% sigma; the laws this clause classifies have none (the Student-t law has no exponential moment
% at any sigma =/= 0), and the reading that does apply to them is the same formula at
% sigma = i omega, whose +-theta pairing the entry says in as many words that it does NOT carry
% -- that pairing is prop:thorin-subclass (1) <=> (2) and is proved there. So a Lean definition
% of the class by A17's formula would be empty at the members, and a definition post-pairing
% would make this lemma a tautology, which is what the withdrawn Skeleton.thorin_subclass_ggc
% was. Should the clause ever need a declaration, the missing step is a formal statement of
% Bondesson's definition at imaginary argument: a librarian read of section 7.1 and a new ledger
% entry, not a Lean exercise. Nothing downstream consumes this lemma.
% CHECK (draft, checks for \S10, item 4): the \Levy-density form of the extended generalized
% gamma convolutions -- that |x| times the two-sided \Levy{} density is completely monotone on
% each half-line -- still wants its page in the source's \S7.1. Neither this lemma nor
% prop:thorin-subclass uses that form; clause (1) there is the folded statement of it.
% NEW as a node (2026-09-15, module B step 2, R160): the elementary half of lem:thorin-ggc, split
% out and machine-checked on Lean core. What ledger A17 says in as many words that it does NOT
% carry -- the pairing of +-theta at sigma = i omega -- is exactly this node. Three findings from
% writing it down. The product of the two factors is the POSITIVE REAL (1 + omega^2/theta^2)^{-1},
% so no branch of the complex logarithm has to be chosen and the printed proof's sum of two
% complex logarithms is avoided. The modulus of a SINGLE factor already gives half the folded
% integrand at every theta =/= 0, which is where the folding convention's factor 2 comes from and
% why the symmetry of V is needed only to cancel the imaginary parts. And the fold is stated at
% the measure: U = 2 V|_{(0,infty)}, so that eq:thorin's own factor 1/2 absorbs it.
\begin{lemma}[Thorin's two-sided exponent at $\sigma = i\omega$]\label{lem:thorin-two-sided}
\uses{prop:thorin-subclass}
\lean{SpatialLine.thorin_pair_mul, SpatialLine.neg_log_norm_thorin_factor,
SpatialLine.thorin_pair_compensator, SpatialLine.thorin_two_sided_fold}\leanok
\notes{gamma-matern-corner}
Let $V \ge 0$ be a measure on $\RR$ with $V\{0\} = 0$, invariant under $\theta \mapsto -\theta$,
and let $U := 2\,V|_{(0,\infty)}$ be its fold onto the positive half-line. For $\theta \ne 0$ and
$\omega \in \RR$,
\[
  \frac{\theta}{\theta - i\omega}\cdot\frac{-\theta}{-\theta - i\omega}
   = \Bigl(1 + \frac{\omega^2}{\theta^2}\Bigr)^{-1},
  \qquad
  -\log\Bigl|\frac{\theta}{\theta - i\omega}\Bigr|
   = \frac12\log\Bigl(1 + \frac{\omega^2}{\theta^2}\Bigr),
  \qquad
  \frac{i\omega\,\theta}{1+\theta^2} + \frac{i\omega\,(-\theta)}{1+(-\theta)^2} = 0 ,
\]
and consequently
\[
  -\int_{\RR}\log\Bigl|\frac{\theta}{\theta - i\omega}\Bigr|\,V(d\theta)
   = \frac12\int_{(0,\infty)}\log\Bigl(1 + \frac{\omega^2}{\theta^2}\Bigr)\,U(d\theta) .
\]
So Thorin's two-sided exponent, read at $\sigma = i\omega$ with $b = 0$ and $V$ symmetric, is the
right-hand side of \ref{eq:thorin} at $(a, U)$ with $a = \tfrac12c$: the first identity pairs
$\pm\theta$ without choosing a branch of the logarithm, the third cancels the compensating terms,
and the second says that the real part alone already supplies the folded integrand.
\statusT\quad\emph{(Machine-checked to Lean core. The three pointwise identities are
\texttt{thorin\_pair\_mul}, \texttt{neg\_log\_norm\_thorin\_factor} and
\texttt{thorin\_pair\_compensator}; the fold is \texttt{thorin\_two\_sided\_fold}, which bundles
the Gaussian term with the integral and concludes at \texttt{thorinExponentL}, this development's
name for the right-hand side of \ref{eq:thorin}. Nothing is spent: no ledger entry is on the path.
What this node deliberately does \emph{not} say is that the two-sided form is membership in
Bondesson's class --- that is the reading of a definition, it is \ledger{A17}, and it is
\ref{lem:thorin-ggc}.)}
\end{lemma}

% ADDED (2026-09-15, module B drafting, R165): a proof environment. The node went \leanok on
% 2026-09-15 (R160) with its three identities and the fold machine-checked and no printed proof;
% "see Lean" is not a proof (repo CLAUDE.md), and the paper transcribes the proof of record. The
% proof below is the checked route: the three pointwise identities, then the fold as an even
% integrand against a symmetric measure.
\begin{proof}\leanok
The three identities are pointwise computations at a fixed $\theta \ne 0$ and $\omega \in \RR$.
For the first, the two denominators multiply to $(\theta - i\omega)(-\theta - i\omega)
= -(\theta - i\omega)(\theta + i\omega) = -(\theta^2 + \omega^2)$, and the two numerators to
$-\theta^2$, so the product is $\theta^2/(\theta^2 + \omega^2) = (1 + \omega^2/\theta^2)^{-1}$,
a positive real number: no branch of a logarithm is chosen because none is taken. For the
second, $|\theta - i\omega| = \sqrt{\theta^2 + \omega^2}$, so
$-\log|\theta/(\theta - i\omega)| = \tfrac12\log\bigl((\theta^2+\omega^2)/\theta^2\bigr)
= \tfrac12\log(1 + \omega^2/\theta^2)$. The third is $\theta + (-\theta) = 0$ over the common
denominator $1 + \theta^2$, which is even in $\theta$.

For the fold, the integrand of the left-hand side is, by the second identity,
$\tfrac12\log(1 + \omega^2/\theta^2)$, an even function of $\theta$ that is finite off the
origin; $V$ gives the origin no mass and is invariant under $\theta \mapsto -\theta$, so the
integral over $\RR$ is twice the integral over $(0,\infty)$, and twice $V|_{(0,\infty)}$ is $U$.
That is the displayed identity. Read at Thorin's exponent with $b = 0$, the compensating terms
cancel in the pairs $\pm\theta$ by the third identity, the Gaussian term $\tfrac12 c\,\sigma^2$
at $\sigma = i\omega$ is $-\tfrac12 c\,\omega^2$, and the pairs of logarithms contribute
$-\log(1 + \omega^2/\theta^2)$ each by the first identity, which is twice the modulus identity;
so the \Levy{} exponent, minus Thorin's, is $\tfrac12 c\,\omega^2 + \tfrac12\int_{(0,\infty)}
\log(1 + \omega^2/\theta^2)\,U(d\theta)$, the right-hand side of \ref{eq:thorin} at
$(\tfrac12 c, U)$.
\end{proof}

% CHANGED (2026-09-15, module B step 2, R160): THE NODE IS NARROWED TO THE DEFINITIONAL
% TRANSLATION. It asserted the equivalence outright, [A] on A17, with the pairing computation
% carried in its proof as an unchecked [T] step; the pairing is now lem:thorin-two-sided above,
% machine-checked on Lean core, and what is left here is the one thing A17 supplies and no Lean
% statement can: that membership in Bondesson's class IS the two-sided representation, read at
% imaginary argument. The node stays notready and has no Lean, for the reason the skeleton note
% above gives (a definition by A17's formula would be empty at the members of the family, and a
% definition taken after the pairing would make the lemma a tautology -- the withdrawn
% Skeleton.thorin_subclass_ggc, review R19); A17 stays unadmitted, grounding a name and not a
% fact, and nothing downstream consumes the node. (P-0010: this is the second of the two entries
% it asked about, and the answer is the same -- no admission is needed.)
\begin{lemma}[the Thorin subclass is the symmetric generalized gamma convolutions]
\label{lem:thorin-ggc}
\uses{prop:thorin-subclass,lem:thorin-two-sided}
\notready
\notes{gamma-matern-corner}
A symmetric law is an extended generalized gamma convolution in Thorin's sense exactly when its
exponent has the two-sided form of \ref{lem:thorin-two-sided} --- that is, when it is
$-\tfrac12c\,\sigma^2 - \int_{\RR\setminus\{0\}}\bigl(\log\tfrac{\theta}{\theta-\sigma}
- \tfrac{\sigma\theta}{1+\theta^2}\bigr)V(d\theta)$ with $c \ge 0$ and $V \ge 0$ symmetric, read
at $\sigma = i\omega$. The two minus signs are the convention used here, in which the exponent is
minus the logarithm of the transform; the source writes the logarithm of the moment generating
function, without them. Consequently, for an admissible $F$ with data $(a,k)$, the two equivalent
conditions of \ref{prop:thorin-subclass} hold if and only if the kernels of the family are
symmetric extended generalized gamma convolutions.
\statusA\quad\ledger{A17}\quad\emph{(\textbf{Assignment.} \ledger{A17} carries the first
sentence and nothing else: the definition of the extended generalized gamma convolutions, the
identification of the symmetric members of that class, and the moment generating function that
defines membership. The passage from that definition to \ref{eq:thorin} --- the pairing of
$\pm\theta$ at $\sigma = i\omega$, which the entry says in as many words that it does not carry
--- is \ref{lem:thorin-two-sided}, machine-checked on Lean core since 2026-09-15; the
equivalence of the two integrability conditions is \textbf{[T]} (2026-09-19, R175: it was recorded
here and in the proof as a term-by-term correspondence, which it is not --- the source's
$\int(1+\theta^2)^{-1}V(d\theta)$ runs over all of $\RR\setminus\{0\}$, where \ref{eq:thorin}'s
condition ran over $(1,\infty)$ alone until the same day).
\emph{Why \texttt{notready}, and why it will stay so.} The entry's statement-as-used is a
\emph{moment generating} function on real $\sigma$, and the laws this lemma classifies have no
exponential moment at any $\sigma \ne 0$, so a Lean definition of the class by that formula would
be empty at the members of the family; a definition taken after the pairing would make the lemma
a tautology, which is what the withdrawn \texttt{Skeleton.thorin\_subclass\_ggc} was (statement
review, R19). What is missing is a formal statement of Bondesson's definition at imaginary
argument, which is a reading of the source and not a Lean exercise. Nothing downstream consumes
this lemma, \ledger{A17} is not admitted, and its name is not on the trust boundary.)}
\end{lemma}

\begin{proof}
By A17 a law belongs to the class exactly when its exponent is
$-b\sigma - \tfrac12 c\sigma^2 - \int\bigl(\log\tfrac{\theta}{\theta-\sigma}
- \tfrac{\sigma\theta}{1+\theta^2}\bigr)V(d\theta)$ for a positive $V$ on $\RR\setminus\{0\}$,
and among these the symmetric ones are those with $b = 0$ and $V$ symmetric. Read at
$\sigma = i\omega$, \ref{lem:thorin-two-sided} turns that into
$\tfrac12 c\,\omega^2 + \tfrac12\int_{(0,\infty)}\log(1 + \omega^2/\theta^2)\,U(d\theta)$ with
$U$ the fold of $V$: the compensating terms cancel in pairs, and the pairing of $\pm\theta$
contributes $-\log(1 + \omega^2/\theta^2)$, which is twice the modulus identity there. That is
\ref{eq:thorin} with $a = \tfrac12 c$.

The source's integrability condition on $V$ and \ref{eq:thorin}'s on $U$ are equivalent, and the
comparison is not term by term. The source asks for
$\int_{\RR\setminus\{0\}}(1+\theta^2)^{-1}V(d\theta) < \infty$ and
$\int_{|\theta|\le1}|\log\theta^2|\,V(d\theta) < \infty$; folded, these read
$\int_{(0,\infty)}(1+\theta^2)^{-1}U(d\theta) < \infty$ and
$\int_{(0,1]}2\log(1/\theta)\,U(d\theta) < \infty$. Each direction is an elementary comparison of
integrands. Forwards: on $(0,1]$ one has $\log(1+\theta^{-2}) \le \log 2 + 2\log(1/\theta)$ and
$(1+\theta^2)^{-1} \ge \tfrac12$, so the first condition bounds $U((0,1])$ and the second the
logarithm; on $(1,\infty)$, $\log(1+\theta^{-2}) \le \theta^{-2} \le 2(1+\theta^2)^{-1}$.
Backwards: $\log(1+\theta^{-2}) \ge \log 2$ on $(0,1]$ and
$\ge \tfrac12\theta^{-2} \ge \tfrac12(1+\theta^2)^{-1}$ on $(1,\infty)$, while
$2\log(1/\theta) \le \log(1+\theta^{-2})$ on $(0,1]$. So membership in the class is
\ref{prop:thorin-subclass}(2), which is (1) by that node.
\end{proof}

The bijection in clause (3) is why the causal corners land on named spatial families with the
same organisation: the causal Gamma family is the single atom of the causal Thorin subclass and
the \Matern{} family the single atom of the symmetric one; the stable families correspond under
$\theta \mapsto \sqrt{2\theta}$; and the causal Bessel family, \emph{if} its delay law is a
generalized gamma convolution (\ref{rem:student-ggc}), maps to a Student-t family that is Thorin
too. What clause (4) adds is that the bridge is a bijection of the
\emph{Thorin} subclasses and an injection, not a bijection, of the whole cones.

% draft: Proposition 10.5
% twin: hcs prop:bessel-family (ledger A8 there)
% CHANGED (review 2026-09-09, R25): Skeleton.bridge_families_bessel is added to the \lean tag.
% Clause (1)'s identification of the kernel at canonical scale 1 IS that declaration -- the
% conditioning computation, shared with prop:bridge-families(3) -- while student_density carries
% only the scaling phi_t = t^{-1}phi_1(./t). Without it the tag covered clause (1) by halves.
% CHANGED (proof 2026-09-10, wave 7): Skeleton.student_transform becomes
% SpatialLine.student_transform -- clause (2) is proved -- and the annotation records which
% ledger entries are off its path. The proved declaration's TYPE is not the skeleton's: the
% causally admissible F and its Laplace-transform specification are dropped, neither being
% mentioned by the conclusion. That is a strengthening (a hypothesis removed), so nothing that
% cited the clause weakens, and the statement of record above is untouched. The node stays
% notready on clauses (3) and (4) and on the causal input Skeleton.student_causal.
% CHANGED (2026-09-19, article mirror of module B, R174): clause (4)'s closing commentary -- "and
% the shape parameter is a dial on the number of finite moments, twice the causal count" -- is moved
% into the annotation (the register pass's item 13), where the comparison with the causal count is
% stated with both inequalities. The clause keeps the moment dichotomy and the tail, which is what
% Skeleton.student_moments is typed on. Safe direction: the statement says less.
% CHANGED (2026-09-15, module B step 2, R160): CLAUSE (3) IS CONDITIONAL ON THE THORIN
% REPRESENTATION OF THE DELAY LAW, and the node leaves ledger A18 for \statusT. The clause
% asserted self-decomposability and membership in the symmetric Thorin subclass outright, on the
% strength of A18; it is now stated under the hypothesis that the inverse-gamma delay law is a
% generalized gamma convolution -- which is what A18 says, and is now cited at the remark
% rem:student-ggc below rather than admitted -- and is machine-checked under that hypothesis as
% SpatialLine.student_thorin, spending only A3 and A11 through thorin_bridge, both already on the
% trust boundary. Skeleton.student_causal is WITHDRAWN: it asserted the existence of a causally
% admissible exponent with the inverse-gamma Laplace transform, which is self-decomposability
% alone, while this clause's Thorin half needs the causal profile to be a Laplace transform; the
% interface as typed dropped exactly what the clause consumes, and no route from it to the clause
% existed. The node stays notready on clause (4) alone (Skeleton.student_moments).
% CHANGED 2026-09-19 (R197; the author's decision D7, the notation half): ONE SENTENCE IS ADDED
% at the end of the statement, saying which of the two things the letter a names here. The
% external review's notation list has "a is the Gaussian coefficient everywhere and the Student-t
% shape in Lemma 4.8, Proposition 5.10, Remark 5.11 and Table 3, and both live in the same
% subsection". The author's decision is to keep the letter -- it is the causal article's and the
% literature's for the shape -- and to say so here. The sentence claims nothing new: that the
% Gaussian coefficient vanishes is what the proof of clause (3) computes under that clause's
% hypothesis, and the sentence says so rather than asserting it unconditionally.
\begin{proposition}[the Student-t family]\label{prop:student-t}
\uses{prop:bridge-families,lem:bridge-exponents,prop:thorin-subclass,prop:moments-tails}
\lean{SpatialLine.bridge_families_bessel, SpatialLine.student_density,
SpatialLine.student_transform, SpatialLine.student_thorin, Skeleton.student_moments}\notready
\notes{bessel-student-corner}
For $a > 0$ let $T_1$ be the scaled inverse-gamma law of the causal Bessel family, with density
$p_{T_1}(u) = u^{-a-1}e^{-1/(2u)}/(2^a\Gamma(a))$, and let $X_1 := \BM_{T_1}$. Then:
\begin{enumerate}
\item the kernel at canonical scale $1$ is
$\phi_1(x) = \frac{\Gamma(a+\frac12)}{\sqrt\pi\,\Gamma(a)}\,(1+x^2)^{-a-\frac12}$, the
Student-t law with $2a$ degrees of freedom in the scale in which the density is proportional to
$(1+x^2)^{-a-1/2}$, and $\phi_t = t^{-1}\phi_1(\cdot/t)$;
\item the transform is
$e^{-F(\omega)} = \frac{2^{1-a}}{\Gamma(a)}\,|\omega|^aK_a(|\omega|)$;
\item if $T_1$ is a generalized gamma convolution --- \ref{rem:student-ggc} --- then the law is
self-decomposable and lies in the symmetric Thorin subclass;
\item $\mathbb{E}|X_t|^n < \infty$ if and only if $n < 2a$, and the tail is
$\phi_t(x) \sim c\,t^{2a}|x|^{-2a-1}$, polynomial.
\end{enumerate}
Throughout, $a$ is the shape parameter of the causal Bessel family --- the letter that family
carries in the literature --- and not the Gaussian coefficient of \ref{eq:sd-profile}, which is a
different number and vanishes for these members, as the proof of clause (3) computes under that
clause's hypothesis.
\statusT\quad\emph{(No ledger entry is spent, and the one fact this article does not prove about
the corner --- that the inverse-gamma delay law is a generalized gamma convolution, \ledger{A18}
--- is the \emph{hypothesis} of clause (3), cited at \ref{rem:student-ggc} and admitted nowhere.
Clause (1) is the conditioning computation written out below; clause (2) is the causal transform
at $\sigma = \omega^2/2$; clause (3) is \ref{lem:bridge-exponents} for the
self-decomposability and \ref{prop:thorin-subclass}(3) for the Thorin half; clause (4) is the
tail of clause (1) with \ref{prop:moments-tails}(1). The shape parameter is therefore a dial on
the number of finite moments, and the count is twice the causal one, $n < 2a$ against $n < a$ for
the delay; that sentence stood inside clause (4) until 2026-09-19 (R174).
\emph{What is machine-checked}: clause (1) entire, as
\texttt{SpatialLine.bridge\_families\_bessel} with \texttt{SpatialLine.student\_density}, and
\textbf{clause (2)}, as \texttt{SpatialLine.student\_transform}, both on Lean core (2026-09-10);
and \textbf{clause (3)} under its hypothesis, as \texttt{SpatialLine.student\_thorin}
(2026-09-15), which spends \ledger{A3} and \ledger{A11} through
\texttt{thorin\_bridge} and no name of its own. Neither
\ledger{A18} nor \ledger{A16} is on the path of clause (2): the transform is obtained from the
conditioning identity of clause (1) and the Laplace transform of the inverse-gamma delay law,
and that Laplace transform is, after the exponential change of variables, the integral by which
$K_a$ is \emph{defined} in this development --- so no asymptotic of the Bessel function is
needed, and the causal exponent is never formed. \emph{The formal statement of (3) says two
things more than the display:} the Gaussian coefficient of the spatial member vanishes, and its
Thorin measure is twice the image of the causal one under $\theta \mapsto \sqrt{2\theta}$, both
of which the checked route produces --- \ref{prop:thorin-subclass}(3) read at the hypothesis.
The node stays \texttt{notready} on clause (4) alone.)}
% CHECK (draft, checks for \S10, item 7): the scaling convention, so that the degrees of
% freedom come out as $2a$ and not $a$, is fixed by the density in clause (1) and the rate in
% $p_{T_1}$; the two should be read together whenever the family is quoted elsewhere.
\end{proposition}

\begin{proof}[Proof of the parts held \textbf{[T]}]
(1) Conditioning on $T_1$, $\phi_1(x) = \int_0^\infty g_u(x)\,p_{T_1}(u)\,du$; the substitution
$w = 1/u$ turns this into
\[
  \frac{1}{\sqrt{2\pi}\,2^a\Gamma(a)}\int_0^\infty w^{a-\frac12}e^{-w(1+x^2)/2}\,dw
  = \frac{\Gamma(a+\frac12)}{\sqrt{2\pi}\,2^a\Gamma(a)}\Bigl(\frac{2}{1+x^2}\Bigr)^{a+\frac12}
  = \frac{\Gamma(a+\frac12)}{\sqrt\pi\,\Gamma(a)}\,(1+x^2)^{-a-\frac12}.
\]
The scaling clause is \ref{prop:bridge-families}. (2) is the causal transform
$e^{-F_{\mathrm I}(\sigma)} = \frac{2^{1-a}}{\Gamma(a)}(\sqrt{2\sigma})^aK_a(\sqrt{2\sigma})$
evaluated at $\sigma = \omega^2/2$; the checked proof computes that Laplace integral rather than
quoting it, and the computation is one substitution: with $A = 1/2$, $B = \sigma$ and
$u = c\,e^{v}$, $c = \sqrt{A/B}$, the integral
$\int_0^\infty u^{-a-1}e^{-1/(2u) - \sigma u}\,du$ becomes
$2c^{-a}\int_0^\infty\cosh(av)\,e^{-2\sqrt{AB}\cosh v}\,dv$, which is the integral by which
$K_a$ is defined here. No asymptotic of $K_a$ enters, at this or any order.

\emph{(3).} Let $U_{\mathrm I}$ be a Thorin measure of $T_1$, as the hypothesis provides:
$\mathbb{E}e^{-\sigma T_1}
= \exp\{-\int_{(0,\infty)}\log(1 + \sigma/\theta)\,U_{\mathrm I}(d\theta)\}$ for $\sigma \ge 0$,
with $\int_{(0,\infty)}\log(1 + \theta^{-1})\,U_{\mathrm I}(d\theta) < \infty$. That exponent is
causally admissible with drift $0$ and delay profile
$k_{\mathrm I}(u) = \int e^{-\theta u}\,U_{\mathrm I}(d\theta)$: a Laplace transform is
nonnegative and nonincreasing, and the causal Frullani identity
$\int_0^\infty(1 - e^{-\sigma u})e^{-\theta u}\,du/u = \log(1 + \sigma/\theta)$ with Tonelli
turns the exponent of \ref{def:causal-admissible} into the Thorin integral, so that the two
integrability windows follow from its finiteness at $\sigma = 1$ --- for $1 \wedge u \le
e\,(1 - e^{-u})$ bounds the windows by the exponent there. By \ref{lem:bridge-exponents} the
image $F_{\mathrm I}(\omega^2/2)$ is admissible, and by the mixture identity with clause (1) it
is an exponent of the Student-t family: the law is self-decomposable. Complete monotonicity of
the spatial profile --- membership in the symmetric Thorin subclass by
\ref{prop:thorin-subclass}(1) iff (2) --- is \ref{prop:thorin-subclass}(3) read at that datum,
which also returns the Thorin measure, twice the image of $U_{\mathrm I}$ under
$\theta \mapsto \sqrt{2\theta}$, and the vanishing Gaussian coefficient.

\emph{(4).} The
tail of the density in (1) is $|x|^{-2a-1}$, so $\mathbb{E}|X|^n < \infty$ exactly when
$n < 2a$, which also follows from $\mathbb{E}|\BM_T|^{2n} = (2n-1)!!\;\mathbb{E}T^n$ and
$\mathbb{E}T_1^n < \infty$ exactly when $a > n$.
\end{proof}

% NEW as a node (2026-09-15, module B step 2, R160): the cited existence statement that
% lem:student-subordinated and prop:student-t(3) now quantify over. It is where ledger A18 binds,
% and it is a REMARK, on the model of rem:bridge-subordination -- cited, marked as not
% machine-checked, consumed by no declaration -- so that the entry grounds no statement node and
% is not admitted. The page anchor and the parameter reading are Halgreen's own: section 2,
% pp. 14-15, at lambda = -a < 0, chi = 1, psi = 0, formulas (2), (3) and (5).
\begin{remark}[the inverse-gamma delay law is a generalized gamma convolution]
\label{rem:student-ggc}
The hypothesis that \ref{lem:student-subordinated} and clause (3) above quantify over is a
theorem of the literature, and it is the one thing about this corner that the article does not
prove. For $a > 0$ the inverse-gamma law $T_1$ of shape $a$ is a generalized gamma convolution:
there is a positive measure $U_{\mathrm I}$ on $(0,\infty)$ with
$\int_{(0,\infty)}\log(1 + \theta^{-1})\,U_{\mathrm I}(d\theta) < \infty$ and
\begin{equation}\label{eq:student-ggc}
  \mathbb{E}e^{-\sigma T_1}
   = \exp\Bigl\{-\int_{(0,\infty)}\log\Bigl(1 + \frac{\sigma}{\theta}\Bigr)
       U_{\mathrm I}(d\theta)\Bigr\}, \qquad \sigma \ge 0 .
\end{equation}
That is ledger A18 --- Halgreen, \S2, pp. 14--15, read at $\lambda = -a < 0$, $\chi = 1$,
$\psi = 0$, the parameters at which the generalized inverse Gaussian density of his (1) is the
density $p_{T_1}$ of \ref{prop:student-t}. His (5) exhibits the measure: it has density
$u(y) = g_a(2y)$ on $(0,\infty)$, where
$g_\nu(x) = 2\{\pi^2x(J_\nu^2(\sqrt x) + Y_\nu^2(\sqrt x))\}^{-1}$ is Grosswald's, and his
integrability conditions (3) --- $\int_0^1|\log y|\,U_{\mathrm I}(dy) < \infty$ and
$\int_1^\infty y^{-1}U_{\mathrm I}(dy) < \infty$ --- are the displayed finiteness for a measure
finite on the compact subsets of $(0,\infty)$; the displayed form is what finiteness of the
exponent at $\sigma = 1$ says, and it is the form both consumers read. Being a generalized gamma
convolution, $T_1$ is in particular self-decomposable, which is the weaker hypothesis
\ref{lem:student-subordinated} quantifies over.

This remark is \emph{not} machine-checked, and it is the only part of the Student-t corner that
is not: the measure Halgreen exhibits is built from the Bessel functions $J_\nu$ and $Y_\nu$,
of which this development has neither, and no route to \ref{eq:student-ggc} through the
elementary theory is known to the author. Nothing in the article depends on the remark ---
both consumers quantify over \ref{eq:student-ggc} or over its consequence instead --- so ledger
A18 is not admitted as a Lean axiom and its name is not on the trust boundary.
\end{remark}

% CHANGED 2026-09-19 (R195, the author's decision at the second review round): the three
% references into chapter 12 are replaced by prose. Module C is not published, and a reference
% from a module-B chapter into a module-C statement would break at a repository split. What the
% paragraph announces is named as further work and is asserted by no statement here.
% TODO(module C): cite thm:scale-locality and thm:joint-locality here once module C is released.
What cuts the Student-t family out of the cone is a locality principle, but not the one the
causal article used for its ancestor. Two principles are in play, and both are left for further
work: locality of the scale evolution, which selects the Gaussian ray alone, and locality of the
scale \emph{space}, as the solution of a second-order equation in $(t,x)$ jointly, which selects
exactly the Gaussian and the Student-t families. Neither is stated here.

% draft: Proposition 10.6
% twin: hcs prop:stable-family + prop:stable-moments
% (Hemigroup.SelfDecomposableExponent.stableExponent_toRealExponent,
%  ...stableExponent_lintegral_pow_kernel)
\begin{proposition}[the symmetric stable family]\label{prop:stable-family}
\uses{cor:semigroup-case,prop:moments-tails,prop:thorin-subclass,prop:bridge-families}
\lean{SpatialLine.semigroup_case_profile, SpatialLine.semigroup_case_gaussian,
SpatialLine.stable_family_moments, SpatialLine.stable_family_thorin,
SpatialLine.stable_family_bridge}\leanok
\notes{covariance-strata}
For $0 < \alpha < 2$ the exponent $F(\omega) = |\omega|^\alpha$ is admissible with $a = 0$ and
profile $k(x) = C_\alpha^{-1}x^{-\alpha}$, where
$C_\alpha = \int_0^\infty (1-\cos u)\,u^{-1-\alpha}du$; its symmetric Thorin density is
proportional to $\theta^{\alpha-1}d\theta$; the kernels are the symmetric $\alpha$-stable
densities; $\mathbb{E}|X_t|^n < \infty$ if and only if $n < \alpha$; and the family is the only
admissible one-parameter semigroup. It is the image under the bridge of the causal stable family
of index $\alpha/2$. The Gaussian ray closes the family at $\alpha = 2$, with $k \equiv 0$ and
$a = 1$.
\statusT\quad\emph{(Everything is \ref{cor:semigroup-case} and \ref{prop:moments-tails}(1),
except the Thorin density, which is the image of the causal one under
\ref{prop:thorin-subclass}(3). \textbf{The phrase ``the kernels are the symmetric
$\alpha$-stable densities'' is nomenclature, not a clause}: it names the classical family whose
exponent is $|\omega|^\alpha$, and no declaration under the tag equates $\mu_{0,t}$ with a
density formula. What the five declarations prove is the exponent, the profile with its
constant, the Thorin density, the moment dichotomy and the Gaussian endpoint; the moment clause
quantifies over any family with the prescribed cosine transform, which is the same family read
by its transform rather than by a density. Absolute continuity itself is
\ref{prop:kernel-regularity}, and the identification of the density with the classical closed
form is not stated anywhere in this article. The heavy tails are the semigroup's doing: within
the cone the semigroup members are the pure-power profiles, the only ones compatible with
stationarity of the increments in the scale parameter, while every non-power profile with
$\int_1^\infty x\,k < \infty$ has a
finite variance. The extended semigroup class beyond $\alpha = 2$ is not in the cone at all.)}
\end{proposition}

% CHANGED (proof of record 2026-09-09, wave 3): a paragraph is added to the Thorin clause,
% giving the route the formalisation took. The printed route -- the image of the causal
% Thorin density under the bridge -- is kept because it is where the clause comes from and
% because it is the reading that makes the corners correspond; the added paragraph is the
% checked one, and it is shorter, needing neither prop:bridge-families(4) nor
% prop:thorin-subclass(4), and needing NO interface at all -- the profile is EXHIBITED as a
% Laplace transform rather than shown to be one, so Bernstein's theorem is not spent.
% Only the POSITIVITY of C_alpha is used, which is why this clause is not blocked on the
% constant that blocks cor:semigroup-case and stable_family_moments.
% CHANGED (proof of record 2026-09-10, wave 4): a closing paragraph is added saying which of
% the routes the proof gives is the machine-checked one, the node having gone leanok this wave.
% Nothing above it is altered: the moment clause was formalised by exactly the argument already
% printed, and the bridge paragraph stays because it is where the Thorin clause comes from.
\begin{proof}\leanok
The exponent, the profile and the constant are \ref{cor:semigroup-case}, as is the statement
that this is the only admissible one-parameter semigroup. The moment criterion applied to
$k(x) = C_\alpha^{-1}x^{-\alpha}$ gives $\int_1^\infty x^{n-1-\alpha}dx < \infty$ exactly when
$n < \alpha$ (\ref{prop:moments-tails}(1)). The causal stable family of index $\alpha/2$ has
$F_{\mathrm I}(\sigma) = \sigma^{\alpha/2}$, whose image is
$2^{-\alpha/2}|\omega|^{\alpha}$ by \ref{prop:bridge-families}(4), a positive multiple of $F$;
the causal Thorin density is proportional to $\theta^{\alpha/2-1}d\theta$, and its image under
$\theta' = \sqrt{2\theta}$, with the Jacobian, is proportional to $\theta'^{\alpha-1}d\theta'$
by \ref{prop:thorin-subclass}(3).

The Thorin density is also reached directly, without the bridge and without
Bernstein's theorem: for $c > 0$ the Gamma integral
$\int_0^\infty e^{-\theta x}c\theta^{\alpha-1}d\theta = c\Gamma(\alpha)x^{-\alpha}$ exhibits the stable
profile as a Laplace transform at $c = (C_\alpha\Gamma(\alpha))^{-1}$, and the Tonelli of
\ref{prop:thorin-subclass} turns that into \ref{eq:thorin}. Only the positivity of
$C_\alpha$ is used, not its value: the integrand is nonnegative, integrable --- at most
$u^{1-\alpha}/2$ near the origin and at most $2u^{-1-\alpha}$ beyond $1$ --- and nonzero on
$(0,1)$.

Machine-checked: the exponent, the profile and the semigroup clause, which are
\ref{cor:semigroup-case}; the moment criterion applied to the pure power, through
\ref{prop:moments-tails}(1) as a cited interface (ledger \ledger{A13}); the Thorin density by
the direct Gamma-integral route of the preceding paragraph; the bridge identity
$(\omega^2/2)^{\alpha/2} = 2^{-\alpha/2}|\omega|^\alpha$; and the Gaussian endpoint. The image
of the causal Thorin density under $\theta' = \sqrt{2\theta}$ is not the checked route: it is
stated on \ref{prop:bridge-families}(4) and \ref{prop:thorin-subclass}(3), which are both
machine-checked, but the checked proof of the Thorin clause does not pass through them, and the
route is kept because it is the reading that makes the corners correspond.
% CHANGED (2026-09-11, draft sync, R105): "neither of them formalised" was stale. Both nodes are
% \leanok, prop:bridge-families(4) as SpatialLine.bridge_families_stable (BridgeCorners.lean) and
% prop:thorin-subclass(4) as SpatialLine.thorin_bridge (ThorinBridge.lean, wave 6), neither
% sorry and both on the axiom guard. What stays true is that stable_family_thorin takes the direct
% Gamma-integral route, so the bridge route is not the checked one; the sentence now says that,
% which is what the draft already read. No claim changed.
\end{proof}

% CHANGED (2026-09-19, article mirror of module B, R174): the citekeys the paper gave this remark
% on 2026-09-18 (the author's review: each corner placed in its literature) are mirrored here, and
% with them the pointer for the Bessel scale-space whose composition law the remark describes. The
% text is otherwise unchanged.
\begin{remark}[the literature of the \Matern{} corner]\label{rem:matern-literature}
The kernels of \ref{thm:matern} are the \Matern{} covariance functions of spatial statistics,
defined there by the spectral density $(1+\theta^2\omega^2)^{-\gamma}$ or as the Green's
functions of a fractional Helmholtz operator (\sscite{whittle1954stationary};
\sscite{lindgren2011explicit}; \sscite{porcu2024matern}). In image analysis the same kernels appear as
Bessel potentials, and the Bessel scale-space built from them (\sscite{burgeth2005bessel})
composes the family along the
\emph{order} $\gamma$ at fixed range --- an additive semigroup in $\gamma$ that is explicitly
not scale invariant --- with no composition law stated for the range family at fixed order. In
natural-image statistics they appear as Bessel K forms, a statistical model of filter responses
rather than a kernel family (see \sscite{pedersen2005alpha}). In none of these is the family derived from axioms on measurement,
and the hemigroup is what makes the derivation possible: composed along the range the family is
a two-parameter cascade whose increments are not stationary in the scale parameter, so no
semigroup axiomatics can
reach it, which is why the literature composed along the order instead. The probabilistic
identity, the symmetric variance-gamma laws (\sscite{madan1990variance}), is what the
classification actually sees.
\end{remark}
